% =====================================================================
%  Harald Høffding, Philosophie og Theologie. En historisk-kritisk
%  Afhandling. Kjøbenhavn: C. W. Stinck's Forlag, 1866. 46 pp.
%  DIPLOMATIC TRANSCRIPTION.
%
%  SOURCE. KB digitisation, copy DA 1.-2.S 3 8°, barcode 110308002769.
%    Library file: Høffding, Harald/philosophie-theologie.pdf
%    sha256 5aea5f48d28a7b36b4f34ae7e96fe1626bbfa9ff7178260fa8202932f2ca6699
%    (SCANS.tsv key hoeffding/philosophie-og-theologie). 52 PDF pages.
%    Page map: printed p. N = PDF N+4 for pp. 3--46. PDF 5 = title page
%    [1], PDF 6 = blank [2], PDF 51 = blank end-paper, PDF 52 = back board.
%    p. 12 carries no printed folio.
%
%  TYPE. Fraktur body; antiqua (Latin, foreign words, "cand. theol.",
%    page references) → \textit{}; letterspacing → \emph{}. Bare "p."
%    before page numbers is antiqua but left plain (corpus convention).
%    Footnotes are numbered 1), 2) ... and restart on every page, hence
%    \setcounter{footnote}{0} at each \opage. Fraktur qv is kept (qvent,
%    Conseqvents, liqvet); ſs/ß → ss in Danish words, ß kept in German.
%    \opage{N} stands at the exact word where printed page N begins; a
%    page turn is never a paragraph break.
%
%  METHOD (2026-09-23). Transcribed from the page images in five batches
%    (pp. [1], 3--11; 12--20; 21--29; 30--38; 39--46), each by an
%    independent subagent working from 300--900 dpi renders. Each batch
%    was checked against the scan's embedded ABBYY text layer with
%    ocrdiff.py (offset 4): 89 candidates in all, 1 corrected (p. 38
%    "relte" → "rette"), 88 the witness's fault (the layer reads ø as o,
%    æ as cr/ce, J for I). Every page was then read a second time against
%    the images by a different agent, word by word, including emphasis,
%    antiqua and paragraphing. The second reading found 3 word errors in
%    44 pp., all silent normalizations, now corrected:
%      p. 19  "indirekt" → "indrekt" as printed (PRINTED AS IS note)
%      p. 25  "consequent" → "conseqvent"
%      p. 33  "Consequents" → "Conseqvents"
%    and one markup inconsistency (p. 27n dash, 42---44). Emphasis and
%    paragraphing were confirmed in every batch; batch 2 (pp. 12--20)
%    has no letterspacing at all.
%    Word-level accuracy after the second reading: 4 errors found in
%    ~44 pp. before correction (3 by second reading, 1 by ocrdiff);
%    none known outstanding. See ACCURACY-LEDGER.tsv.
%
%  OPEN ITEMS (see the % UNSURE: notes in the text).
%    Title page: publisher's initial "C. W." (W or V).
%    p. 5: comma or semicolon after "Yngre".
%    p. 35: "litteraire" (first e blotted).
%    p. 41: "tvert-imod" (r under a blot).
%    p. 9: "videnstabeligt" is pencil-corrected by a reader in this copy;
%    p. 19: "indrekt" likewise. Printer's errors are kept and marked
%    % PRINTED AS IS:.
% =====================================================================

\documentclass[12pt,a4paper]{book}
\usepackage[utf8]{inputenc}
\usepackage[T1]{fontenc}
\usepackage{libertinus}
\usepackage{libertinust1math}
\usepackage{textalpha}   % Greek (pp. 10, 23)
\usepackage{graphicx}    % for \reflectbox (the old Danish »ɔ:« id-est mark)
\DeclareUnicodeCharacter{0254}{\reflectbox{c}}  % ɔ (U+0254) = reversed c
\usepackage[danish]{babel}
\usepackage[protrusion=true,expansion=true]{microtype}
\usepackage{setspace}
\usepackage[margin=1.2in]{geometry}
\usepackage{enumitem}
\usepackage{fancyhdr}
\setlength{\headheight}{28pt}
\addtolength{\topmargin}{-13.5pt}
\usepackage{hyperref}

\hypersetup{
  pdftitle={Philosophie og Theologie},
  pdfauthor={Harald Høffding},
  hidelinks
}

\renewcommand{\thefootnote}{\arabic{footnote})}

\pagestyle{fancy}
\fancyhf{}
\fancyhead[LE,RO]{\thepage}
\fancyhead[RE]{\textit{Philosophie og Theologie}}
\fancyhead[LO]{\textit{Harald Høffding}}

\setstretch{1.3}

\usepackage{marginnote}
\newcommand{\opage}[1]{\marginnote{\footnotesize\textit{[#1]}}}

\begin{document}
\thispagestyle{empty}

% TITLE PAGE (printed [1], PDF 5)
%   Philosophie og Theologie.
%   [short rule]
%   En historisk-kritisk Afhandling
%   af
%   Harald Høffding,
%   cand. theol.          [antiqua]
%   [double rule]
%   Kjøbenhavn.
%   Bog- og Papiirhandler C. W. Stinck's Forlag.
%   Thieles Bogtrykkeri.  [letterspaced]
%   1866.
% UNSURE: publisher's second initial read as Fraktur W (three-stemmed,
%   unlike the B of "Bog-" on the same line); alternative V.
% Printed [2] (PDF 6) is blank.
\begin{center}
{\Huge Philosophie og Theologie.}\par
\rule{4em}{0.4pt}\par\bigskip
{\large En historisk-kritisk Afhandling}\par\medskip
af\par\medskip
{\large Harald Høffding,}\par
\textit{cand. theol.}\par\bigskip
\rule{6em}{0.8pt}\par\bigskip
Kjøbenhavn.\par
Bog- og Papiirhandler C. W. Stinck's Forlag.\par
{\small Thieles Bogtrykkeri.}\par
1866.
\end{center}
\clearpage

% --- p. 3 ---
% (PDF 7; no folio printed; signature mark "1*" at foot)
\opage{3}\setcounter{footnote}{0}%
Naar en Forhandling angaaende et væsentligt Stridspunkt
atter og atter vender tilbage, naar Discussionen atter og atter
optages igjen, da kan det for en udvortes Betragtning synes
trivielt, saaledes altid at strides om det Samme. Men det
vil ogsaa kun være saa for den udvortes, overfladiske Betragtning.
I Virkeligheden maa Discussionen altid vende tilbage
til samme Kampplads med rigere Udrustning, udstyret
med de nye Midler, den fremskridende Udvikling siden sidst
har bragt frem. Dette finder sin Anvendelse, naar man betragter
de forskjellige Vilkaar og Forudsætninger, under hvilke
Stridighederne mellem Philosophie og Theologie ere blevne
førte til de forskjellige Tider, og særligt, naar vi see hen til
den eiendommelige Maade, hvorpaa Striden har udviklet sig
i \emph{vor} Litteratur.

Hvergang der derfor paa Ny bringes Røre i denne Sag,
er det naturligt at vende Blikket tilbage til Sagens tidligere
Historie, for at see, hvorvidt det nye Indlæg virkelig har forstaaet
og benyttet den ret --- thi det er den første Betingelse
for at komme videre. Men her kan det da stille sig saa forunderligt,
at Striden egentlig ikke længere føres om at afgjøre
Sagen selv, men om at afgjøre, hvem der nu egentlig \emph{har}
seiret i den afgjørende Strid. Det gaaer nemlig ikke i den
ideelle Strid som i Kampene i den legemlige Verden, hvor
Afgjørelsen er til at tage og føle paa.

% --- p. 4 ---
\opage{4}\setcounter{footnote}{0}%
Da jeg læste det nye Indlæg\footnote{„Samvittighed og Videnskab“, af A. C. Larsen, \textit{cand. theol.}} i den ovenfor omtalte
Strid, som, efter hvad jeg har hørt, af Mange ansees for at
have en ikke ringe Betydning, navnlig som Angreb paa den
af Prof. R. Nielsen hævdede Opfattelse af Forholdet mellem
Theologie og Philosophie og mellem Tro og Viden, fik jeg et
Exempel herpaa. Uagtet Forf. synes at beklage, at „siden den
Strid, der i sin Tid førtes om Martensens Dogmatik, er al
Polemik om Troens og Videns Forhold til hinanden aldeles
% PRINTED AS IS: the h of "har" before "han" is a small, raised sort
forstummet“, har han dog ikke gjort, hvad der stod i hans
Magt for paa rette Maade at føre Sagen videre: nemlig
klart at gjøre sig Rede for, hvorledes Sagerne staae. Dette
har skadet hans Arbeide saa meget desto mere, som den historiske
Udviklingsgang temmelig klart viser hen til den Anskuelse, der
maa blive Resultatet af den lange, forviklede Strid. En historisk
Fremstilling af denne Udviklingsgang vil paa eengang
kunne vise os, hvilken den seirende Anskuelse maa være, og
afgive en Kritik over det nævnte Skrift.

\begin{center}\rule{4em}{0.4pt}\end{center}

Enhver Studerende, og da særlig Den, der har siddet
paa Discipelbænken saavel i det philosophiske som i det theologiske
Auditorium, maa føle sig opfordret til at klare sig,
hvad det egentlig er, her strides om, og hvorledes han selv
maa stille sig i denne Sag. Dette er ingenlunde let for den
Unge, der væsentlig maa forholde sig lærende; thi her staaer
Autoritet mod Autoritet, og Den, som nu maaskee seer rokket,
hvad han ansaae for at staae fast som Klippen, bliver let
mismodig og vaklende. Men Eet er han da netop herved
forhindret i: hvad Latinerne kalde: \textit{jurare in verba magistri} ---
han kan da ikke sværge til to Autoriteter paa eengang! Saa
seer han da Nødvendigheden af at tage sit Parti efter egen
bedste Overbeviisning.

% --- p. 5 ---
% (folio 5; the text ends after "denne lille Afhandling.", a short
% rule follows, and the rest of the page is blank)
\opage{5}\setcounter{footnote}{0}%
Hvad der nu gjør, at jeg optræder offentlig med min
Opfattelse af Sagen, istedetfor at beholde den for mig selv,
er, at det synes mig passende, om Striden ogsaa udfægtedes
% UNSURE: comma after „Yngre“ is faint; could be a semicolon
mellem de Yngre, efter at Hovedslaget i sin Tid er slaaet af
Mestrene selv.

I denne Henseende har jeg jo en Forgænger i Forfatteren
af det ovenanførte Skrift, som har givet mig Anledning til
denne lille Afhandling.

\begin{center}\rule{4em}{0.4pt}\end{center}

% --- p. 6 ---
% (PDF 10; no folio printed: text sunk, as on p. 3)
% Page references "p." and their figures are antiqua throughout; left upright.
\opage{6}\setcounter{footnote}{0}%
Førend jeg skrider til min egentlige Opgave: at vise,
hvorledes den Anskuelse, der efter min Overbeviisning er den
eneste holdbare, fremgaaer af Stridens hele Gang, især paa
Grund af den Vending, denne i vor Litteratur har faaet ved
S. Kierkegaards Optræden, --- skal jeg kortelig angive, hvorfor
jeg maa ansee Hr. Larsens Indlæg for at være forfeilet.

Forf. af „Samvittighed og Videnskab“ erkjender naturligviis
en Forskjel mellem Tro og Viden; men han vil ikke
indrømme, at de ere absolut ueensartede. Han anseer det for
umuligt at fastholde Overbeviisningen om den moderne Videns
Realitet og Forvisningen om Christendommens Realitet, uden
ved Hjælp af den videnskabelige Theologie, hvor Tro og Viden
kunne forhandle med hinanden (p. 60).

Men paa denne Maade forsoner han kun de to Modsætninger
ved at sammenblande dem. Han gjør Videnskaben
religieus for at kunne gjøre Religionen videnskabelig.

Naar jeg siger, at Forf. gjør Videnskaben religieus, da
mener jeg derved, at han henviser Problemer, som høre under
Philosophien, til Religionen. --- Hvoraf vide vi, at den udvortes
Sandsning har Realitet? Dette Spørgsmaal kan efter Forf.s
Mening ikke afgjøres „ad \emph{Videns} Vei“. Den udvortes Verden
bliver først virkelig til for mig gjennem \emph{Tro} (p. 13--14), og
denne Tro er principielt Et med den, i hvilken Gud er til
for mig (p. 14). Dette er en Tvetydighed. Philosophien
formaaer virkelig „ad Videns Vei“ at løse Objectiveringsproblemet\footnote{Her kan henvises til „Grundideernes Logik“ I, p. 172--177.}.
Og forsaavidt der kan tales om et philosophisk
% --- p. 7 ---
\opage{7}\setcounter{footnote}{0}%
\textit{sapere aude}, en Tro paa Tænkningens Realitet, saa er
denne Tro principielt forskjellig fra den religieuse, om den
end kan støttes ved denne. Dette kan sees af Forholdet
mellem den dogmatiske Philosophie og den kritiske: „die kritische
Philosophie begreift, was die dogmatische geglaubt hatte“\footnote{Kuno Fischer, Geschichte der neuern Philosophie I, p. 241.},
og hvad den dogmatiske Philosophie havde troet paa, var netop
Erkjendelsens Realitet. --- Hvilket er Forholdet mellem Aanden
og Naturen? Herom kan man efter Forf.s Mening Intet
faae at vide i Videnskaben: „thi Naturvidenskaben har som
saadan intet Naturbegreb; den veed som saadan intet om
Naturens Forhold til det Absolute. De Magter, mellem hvilke
Striden føres, uden at nogensinde Forsoning er mulig, det er
Pantheismen og Christendommen“ (p. 38). Dette er ligefrem
urigtigt. Der oversees nemlig, at der er Noget, som hedder
Naturphilosophie: her haves et Naturbegreb, om ellers Physiken
ikke har noget; og her vises Forholdet mellem Idee og Natur.
Her kan man lære, hvorledes den høieste Aarsag virker i og
gjennem alle de endelige Aarsager. Naturphilosophien vil hverken
troe eller begribe Skabelsesdogmet; men den vil begribe Forholdet
mellem Aand og Natur som et Grundforhold. Vil
man kalde dette Pantheisme, saa er det et tvetydigt Udtryk;
thi Naturphilosophien er ikke Religion men en videnskabelig
Begrundelse af Naturens Phænomener.

Overseer Forf. saaledes\footnote{Man kan her ikke lade være at tænke paa Hegels Ord: „Sie ignoriren die Philosophie, aber nur um in ihren willkürlichen Raisonnements nicht genirt zu werden“. Geschichte der Philosophie I, p. 80.} Philosophien, saa er han
des bedre Ven med Physiken. Hvilken \textit{entente cordiale} er
der ikke for ham mellem „de virkelige Naturlove“ paa den
ene Side og Skabelsesdogmet og Miraklerne paa den anden
Side! I Skabelsesdogmet ligger jo, at Naturen er et virkeligt
„Noget“: saa er der ogsaa „virkelige Love“ og en „virkelig
Naturvidenskab“ (p. 33)! Men hvad mon Naturvidenskaben
% --- p. 8 ---
\opage{8}\setcounter{footnote}{0}%
vil sige, naar man lader Dogmet tale ud: „det er Noget ---
og dog Intet for den Almægtige!“ Jeg troer ikke, det nytter
at forklare Naturforskeren, at nu gaae vi over i en anden
Videnskab, hvor man har denne „anden Maade at see Tingene
paa“. --- Og nu Miraklet! Ja, siger Forf., det „stadfæster
netop Naturloven; sætter man Undtagelsen, maa man ogsaa
sætte Regelen“ (p. 37). Et prægtigt Raisonnement --- især
for Tyve og Røvere, der saaledes kunne indlægge sig Fortjeneste
ved at indskærpe Lovens Gyldighed, om end temmelig
„indirecte“!

Forf. har Ret i, at det ikke er Christendom og Naturvidenskab,
„der skal vælges imellem“; men han har ikke Lov
til at sætte „Theologien“ istedetfor „Christendommen“ (hvilket
han gjør p. 37--38); thi hermed forandres hele Sagen.
Striden mellem Theologie og Naturvidenskab er nemlig en
Strid mellem to Videnskaber.

Men det er endda ikke saa underligt, at Forf. --- bevidst
eller ubevidst --- identificerer Christendom og Theologie. Thi
han har kun gjort Videnskaben religieus for at gjøre Religionen
videnskabelig.

Naar Forf. vil gjøre Religionen til Gjenstand for Videnskaben,
da er det ingenlunde hans Mening, at man „saadan
uden Videre“ kan begribe de religieuse Sandheder. Men han
mener, at „det er noget ganske Andet, naar jeg efter i Troen
at være gaaet ud over mig selv og efter at have sat Gud,
som Aarsag, saa bagefter af min Gudsbevidsthed, som Virkning,
slutter tilbage til Beskaffenheden af Guds Væsen“ (p. 20).
Ved denne megen Gaaen ud og ind opnaaer man saa „en
Videnskab, der væsentlig intet Andet vil være, end selve det
religieuse Livs klare og gjennemsigtige Bevidsthed om sig
selv“ (p. 54).

Men, maa man spørge, gaaer det da lettere med at bevise
og begrunde, efter at man er blevet troende, end før?
Man lægge vel Mærke til, at der ikke spørges: om den
% --- p. 9 ---
\opage{9}\setcounter{footnote}{0}%
religieuse Bevidsthed er overbeviist om Sandheden af Det,
hvorpaa den troer, men: om denne Overbeviisning kan blive
en videnskabelig Indsigt og Begriben?

Nu siger Johannes Climacus, at den christelige Sandhed
er Paradoxet! Denne Hindring for „det religieuse Livs
klare og gjennemsigtige Bevidsthed om sig selv“ maa da først
ryddes afveien. Og det gaaer nemt nok: „Et er det med
lyrisk-lidenskabelig Energi at udtale sig i polemisk Situation;
% PRINTED AS IS: "videnstabeligt" (long-s + t ligature for sk); a reader has
% pencilled a correcting "k" in the margin
et ganske Andet er det at tale positivt-videnstabeligt“ (p. 49).

Heri er vist Ingen uenig med Forf.; men naar han
deraf slutter, at Joh. Climacus's Sætning kun kan gjælde,
naar det er Humoristen, der har den i Munden „i polemisk
Situation“, saa er dette\footnote{Hvorledes mon Forf. har forstaaet „den indirecte Meddelelses“ Væsen og Betydning, naar han kan slutte saaledes!} en forunderlig Blanding af Tankeløshed
og Sophisterie. Den Troendes Forhold til Troens
Gjenstand i Tro og Haab er et ligefremt Forhold, hvori der
% PRINTED AS IS: "Parodox" for Paradox
intet Parodox er; men saasnart den Troende vil have objectiv
(ɔ: videnskabelig)\footnote{Hvad Ret har Forf. til at definere den objective Vished som „reen logisk Tilstandsvished?“} Vished, saa viser Troesgjenstanden sig
--- efter Joh. Climacus --- som det Absurde. Men, siger
Forf., Joh. Climacus siger jo selv, at „den objektive Ligegyldighed“
slet Intet faaer at vide om Christendommen; hvor
kan man saa objektivt faae at vide, at den er det Absurde?
Svaret er endda ikke saa vanskeligt: thi det, at man faaer at
vide, at den er det Absurde, og det, at man slet Intet faaer
at \emph{vide}\footnote{Naturligviis videnskabeligt. --- Exempler hos Joh. Climacus p. 160 f. 166 f.}, er Et og det Samme. Paradoxet maa for den
videnskabelige Domstol erklæres for et \textit{non-ens}. --- Ikke desmindre
betragter Forf. Joh. Climacus som lykkelig og vel
anbragt paa „en egen Plads“ udenfor det „Positiv-Videnskabelige“.
Det er kun ligesom til Overflod, at den stakkels
% --- p. 10 ---
\opage{10}\setcounter{footnote}{0}%
Humorist maa høre, at han ikke har „positiv-videnskabelig“
Alvor nok til at skjelne mellem θεός og ὁ θεός (p. 53); og
saa faaer han til Slutning den Lærdom, at man dog ogsaa
skal have det Objektive med. Og deraf, at Troen ikke kan
undvære „det Objektive“, sluttes\footnote{Paa samme underfundige Maade satte Forf. før Ordet „Theologie“ istedetfor Ordet „Christendom“.} saa ligefrem: at den heller
ikke kan undvære „Videnskabeligheden“ (p. 54)!

Men nok er det: Forf. er blevet færdig med Johannes
Climacus; han etablerer nu Dogmet som en Kjendsgjerning
for den religieuse Bevidsthed, og engagerer Forstanden (efter
Forf.s egen Erklæring p. 52 den samme\footnote{Dette har forresten ikke Noget at sige, da Forf. ved „Forstand“ forstaaer noget aldeles Formelt (p. 50--51). En saa „reduceret“ Forstand kan man faae til, hvad det skal være. Forf. synes ikke at vide, at der er Noget, som hedder philosophisk Tænkning, Begreb, Ideeerkjendelse. Han gjør sig virkelig Opgaven altfor let!} Forstand, som
den Troende havde, før han blev troende) til at tage fat
(p. 50--52; 71--74). Men her maa man beklage, at Forf.
netop har overseet, hvad det kommer an paa. Han stiller
uden Videre Religionen sammen med Naturen, uden at tænke
paa den væsentlige Forskjel mellem begge, betragtede som Gjenstande
for Videnskaben\footnote{Dette kan saameget mindre undskyldes, som denne Parallel er bestridt i en særegen Afhandling: „Er Religionsforskerens Forhold til den saakaldte Aabenbaring det samme som Naturforskerens til Naturen?“ Af Mag. Stilling. 1853.}. Om den religieuse Kjendsgjerning
gjelder, hvad Feuerbach siger: „Thatsache ist kein Grund“.
I Naturens Kjendsgjerninger kan Fornuften derimod finde sit
eget Væsen og saaledes begrunde dem. „Die Vernunft ist
die Hebamme der Natur“. Og som det gaaer med Kategorien
„Kjendsgjerning“, gaaer det med den beslægtede Kategorie
„Forudsætning“. En videnskabelig Forudsætning maa heelt og
holdent kunne gjennemgaae Tvivlens Ildprøve; dette gjælder
% --- p. 11 ---
% (page turn inside "Sandse-/indtryk")
ogsaa om, hvad der er Gjenstand for „de umiddelbare Sandse\opage{11}\setcounter{footnote}{0}indtryk“:
det anerkjendes ikke videnskabeligt, uden det kan begrundes.
Men det kan Forf. ikke gaae ind paa: „hvor berettiget
man end er til at forlange Tvivlen, Et maa man dog
ikke forlange, at den religieuse Bevidsthed skal tvivle om Realiteten
af sin umiddelbare Væren“ (p. 58.). Jo, det maa og
skal man forlange, hvis der skal komme Videnskab ud af det,
og i denne Henseende afviger Theologien saare meget i sine
„Fordringer“ fra Naturvidenskaben.

En praktisk Prøve paa Forf.s Videnskab faae vi i hans
3die Afhandling, hvor han vil vise os, hvorledes Gud staaer
„i realt Vexelforhold“ med Tiden (p. 46). Kan han nu
virkelig udvikle dette Vexelforhold mellem den evige, personlige,
alvidende Gud og den timelige, successive Udvikling, saaledes
at det bliver et videnskabeligt Begreb? Han synes dog kun at
vise, at det „maa tænkes“ --- men han tænker det ikke. ---

Naar Forf. saaledes hverken har givet Videnskaben eller
Troen, hvad dem hver for sig tilkommer, kan Forsoningen
heller ikke blive varig. Og at dette er saa, vide vi ikke
blot \textit{a priori}, men --- som det Følgende viser --- ogsaa
\textit{a posteriori}.

% p. 11 ends with a short rule; the rest of the page is blank.
\begin{center}\rule{4em}{0.4pt}\end{center}

% ==== batch 2: printed pp. 12--20 (PDF 16--24) ====
% p. 11 closes its section with a short rule; p. 12 opens a new section
% with a sunk start and carries no printed folio (only set-off from p. 11
% shows at its head). No heading is printed.
% --- p. 12 ---
\opage{12}\setcounter{footnote}{0}%
I Middelalderen var Theologien som Scolastik den herskende, ja den eneste
Videnskab. „Dogmet er sandt (thi Kirken lærer det!); altsaa er det ogsaa
stemmende med Fornuften; altsaa kan det bevises“. Saaledes sluttede man, og
opbød nu alle Forstandens Kræfter for at opbygge de store scolastiske
Systemer, man træffende har sammenlignet med de gothiske Domkirker. Og det
var en from Gjerning; ligesom de fromme Bygmestre ved Rhinen reiste de mod
Himlen pegende Spiir, hvor der ofte høit oppe, hvor aldrig noget Menneskes
Øie naaede, var anbragt Prydelser til Guds Ære, saaledes opførtes ogsaa de
scolastiske Systemer, med deres hele Apparat af Syllogismer og
Distinctioner, \textit{in gloriam Dei}. Scolastikerne kæmpede med de
gjenstridige Tanker paa samme Viis som Korsfarerne mod Saracenerne, og ilde
gik det begge, hvis de ikke bøiede sig for Troen og Kirken.

Med den middelalderlige Kirke faldt ogsaa Scolastiken. Dens Fald bevirkedes
baade af indre Selvopløsning og ydre Angreb. Scolastikens Forudsætning var
Eenheden af Tro og Viden: men nu havde allerede Duns Scotus ved, i
Modsætning til Thomas Aqvinas, at gjøre Villien til høieste Princip
skjelnet mellem det Theoretiske og det Practiske, og Vilhelm Durandus (i
første Halvdeel af 14de Aarhundrede) havde erkjendt, at der gaves Dogmer,
der ikke kunde og ikke skulde bevises, men troes som aabenbarede Sandheder.
Scolastikens Forudsætning var endvidere Eenheden af Tænken og Væren:
% --- p. 13 ---
\opage{13}\setcounter{footnote}{0}%
men nu var ved Vilhelm Occam Nominalismen (\textit{universalia sunt nomina,
flatus vocis}) bleven herskende i skærpet og fornyet Skikkelse. Angrebene
udenfra skete deels fra den friblevne Tankes Side, som nu vilde nedrive det
Fængsel, hvor den saalænge var bleven martret, og deels fra det
existentielle, religieuse Standpunkt, idet Luther brød med det traditionelle
Dogma og alle de »\textit{traditiones humanæ}«, hvori det havde spundet sig
ind.

Men den nye philosophiske Forsken vovede endnu ikke at angribe Kirkelæren
selv. For da at kunne fastholde deres egne Erkjendelsesresultater uantastede
af Kirkens Myndighed, opstillede Philosopherne Adskillelsen mellem
\textit{duæ veritates, philosophica et theologica}, ved den Sætning, at
Noget kunde være sandt i Theologien, fra Dogmets Standpunkt, og dog falskt,
eller idetmindste ubevisligt fra Philosophiens, Fornuftens Standpunkt. Som
Mænd, der hævdede denne Adskillelse, kan nævnes Pietro Pomponazzo (i
Slutningen af 15de og Begyndelsen af 16de Aarhundrede), der især
beskjæftigede sig med Udødelighedslæren, og Vanini (i Slutningen af 16de og
Begyndelsen af det 17de Aarhundrede), hvem den omtalte Distinction ikke
reddede fra Baalet.

At denne Distinction var sophistisk og uholdbar, er let at see. Thi Dogmet,
saaledes som Kirken forelagde det til lydig Antagelse for sine Troende, var
ikke den evangeliske Sandhed med dens practisk-religieuse Characteer, men det
var et Theologumenon med kirkelig Sanction. „Den videnskabelige og den
religieuse Erkjendelse vare den Gang saa kirkelig sammenvoxne, at de ikke
kunde skilles ad. Naar Philosopherne da, for at bevare Videnskabeligheden,
erklærede, at de i alle Troessager vare villige til at underkaste sig
Kirkens Myndighed, saa var man fra Kirkens Side vistnok beføiet til at drage
Oprigtigheden af en saadan Indrømmelse i Tvivl: hvorledes skulde Nogen vel
kunne holde fast ved en Fornuftsandhed og dog anerkjende en theologisk
Sandhed, naar de staae i af%
% --- p. 14 ---
\opage{14}\setcounter{footnote}{0}%
gjørende Strid med hinanden, og dog begge skulde gjælde for
videnskabelige?“\footnote{R. Nielsen: Grundideernes Logik p. XXII.}

At Striden paa disse Vilkaar ikke kunde føre til Forsoning, er da
indlysende. Den fortsattes da saavel i den protestantiske Kirke som i den
katholske. Ved Siden af Philosophien stillede sig nu Naturvidenskaben, som
først fra det 16de Aarhundrede af har en fortløbende Historie.

Hvis der i Slutningen af forrige Aarhundrede var blevet opkastet det
Spørgsmaal, hvem der havde seiret i Striden, Theologien eller Philosophien,
kunde Svaret neppe være tvivlsomt. Theologien havde nemlig, under
Indflydelse af Tidsaanden, ladet det ene Mirakel falde efter det andet,
saavel i Aandens som i Naturens Verden. Det philosophiske, naturalistiske
18de Aarhundredes Theologie var ikke længere nogen Theologie.

Men det almindelige aandelige Opsving i Begyndelsen af dette Aarhundrede
viste sig ogsaa paa Theologiens Omraade. Det var især Philosophien, der nu
blev af Betydning. Man havde forkastet Rationalismen og atter hyldet
Troesprincipet; men man saae i den nyere Philosophie en Mulighed til at
forsone dette Troesprincip med Vidensprincipet. Nu kom der da en
Blomstringstid for Theologien, fuld af Begeistring for dens Dyrkere. „Af de
saakaldte „bedre Stænder“ havde man ingensinde seet saa Mange vælge dette
Studium som i hiin\footnote{Der tales om Aarene 1824--33 i Tydskland.}
Periode.... Ligesom nogle Aar tidligere den tydske Ungdoms Blomst havde
stillet sig under Krigens Faner for at frelse Fædrelandet fra politisk
Trældom, saaledes ansaae i de derpaa følgende Fredsaar de rigest begavede i
aandelig Henseende det for den høieste Ære og høieste Vinding at turde tjene
Kirken og offre den sine bedste Kræfter.... Man vidste, at Kampen gjaldt et
ædelt Maal; man troede paa Muligheden af en endelig Seir, paa Muligheden af
en Forsoning af Tro og
% --- p. 15 ---
\opage{15}\setcounter{footnote}{0}%
Viden, om man end ikke saae denne fuldfærdig;... man ansaae det endnu ikke
for at bære et fremmed Aag, naar Theologien gik Haand i Haand med
Philosophien\footnote{Hagenbach: „Om Aftagelsen af det theol. Studium“, i
Clausens Tidsskrift for udenl. theol. Litt. 1858, p. 13, 21.}.“

Det var især den hegelske Philosophie, som var rig paa Forjættelser om en
Forsoning af Tro og Viden. Dens Indhold var det Absolute --- men ogsaa
Christendommens Indhold er jo den absolute Sandhed\footnote{Det hedder om
den religieuse Aand, at dens Indhold er „als religiös wesentlich
spekulativ“, og om Philosophiens Indhold, at det er „spekulativ und damit
religiös“. Hegel, Encyclopädie der philos. Wissenschaften. Berlin 1845.
p. 513.}. Den havde derhos i sin dialectiske Methode Evnen til at
gjennemføre Mediationen mellem de stridende Principer. Hvor rige
Forudsætninger var der da ikke her for en spekulativ Theologie!

Men satte da Hegel slet ingen Forskjel mellem Tro og Viden? Jo, men en
tvetydig, som derfor ogsaa siden blev fortolket \textit{in utramque partem}.
„Worauf es ganz allein ankommt, ist der Unterschied der Formen des
spekulativen Denkens von den Formen der Vorstellung und des reflectirenden
Verstandes“\footnote{Anf. Skr. p. 512.}. Dette blev nu forstaaet
paa tvende Maader. Nogle sagde: „Forestillingen er en lavere, følgelig en
inadæqvat, altsaa en relativ usand Form, som derfor maa opløses af
Tanken“. Men Andre meente, at man vel kunde beholde den hegelske
Bestemmelse af Troen som en umiddelbar Viden, „kun at vi da ikke ved det
Umiddelbare forstaae det Lavere, det Ufuldkomne, der skal afskaffes af
Philosophien som det Fuldkomne, men den oprindelige, den primitive Viden,
der er Spekulationens Fundament\footnote{Martensen, den christelige
Dogmatik. 1849. p. 13.}.“

Denne sidste Opfattelse blev da Grundlaget for en ny Retning i den
theologiske Dogmatik, som ogsaa fandt Repræsentanter ved vort Universitet.
Det var især nuværende Biskop
% --- p. 16 ---
\opage{16}\setcounter{footnote}{0}%
\textit{Dr.} Martensen, der ved sin\footnote{\textit{De Autonomia
conscientiæ sui humanæ in theologiam dogmaticam nostri temporis introducta.
Hauniæ 1837.}} Afhandling for Licentiatgraden banede Veien hertillands for
denne „nye Æra“. Han vilde ingenlunde tilkjende Philosophien Primatet over
Theologien, men derimod hævde et eget Princip for den theologiske
Spekulation; medens den philosophiske Tænkning er autonom, er den
theologiske theonom (\textit{dogma de deo, auctore conscientiæ sui humanæ,
verum esse theoriæ cognitionis fundamentum, qvo efficitur nexus necessarius
inter scientiam theologicam et auctoritatem dogmatis, vinculum indissolubile
scientiæ et conscientiæ}\footnote{Anf. Skr. p. 135.}.
% PRINTED AS IS: the parenthesis opened before "dogma de deo" is never
% closed; the print has only the footnote mark "2)" before the full stop.
Mediationens Begreb bibeholdtes, men der fordredes, at „Mediationen skulde
søges fra den troende Erkjendelses Standpunkt“\footnote{Dogmatiske
Oplysninger af H. Martensen. 1850. \textit{pag.} 66--69. --- „Christendommens
Middelpunkt, Incarnationslæren, viser jo netop, at den christelige
Metaphysik ikke kan hvile i et Enten---Eller“, saaledes som „den jødiske
Religions Metaphysik.“ Tidsskrift for Litteratur og Kritik 1839, p.
458.}.

Her var altsaa fastsat et dogmatisk Princip, som formeentlig kunde hævde
Theologiens videnskabelige Characteer ved i sig at forene Troens Princip med
Videns. Men hvad vor Litteratur saae endnu i de første Aar af Productioner i
denne Retning, var kun enkelte\footnote{Her kan f. Ex. nævnes „Mester
Ekkardt“, det fortrinlige „Bidrag til at oplyse Middelalderens Mystik,“ af
\textit{Dr.} Martensen (1840).} Forsøg, ikke nogen gjennemført videnskabelig
Fremstilling af Troeslæren efter det angivne Princip. „Systemet var endnu
ikke færdigt“ --- men det ventedes, og Alle imødesaae med store
Forhaabninger den videre Udvikling af den „nye Æra“.

Førend vi fortsætte den historiske Fremstilling videre, maae vi først
erindre om den Retning, der --- ligeledes udspringende fra den hegelske
Philosophie --- gik den modsatte Vei af den „troende“ Spekulation. Det er
Strauß og Feuer%
% --- p. 17 ---
\opage{17}\setcounter{footnote}{0}%
bach, som her ere Hovedmændene. Var det Svar, Viden gav, naar man forlangte
dens Vidnesbyrd, efter at have stillet den overfor Aabenbaringen, for hiin
Retning ubetinget bekræftende, saa blev det her ligesaa ubetinget
benegtende. Den negative Kritik stiller sig altsaa mod den positive
Dogmatik, \textit{contra} mod \textit{pro}.

Modsætningen er da her saa stærk som mulig. Men skulde der dog ikke være
noget Fælles for de tvende stridende Parter, Noget, hvorom de trods al
Uenighed ere enige?

Hvad begge Parter vare enige om, var, at Striden skulde udfægtes paa
Videnskabens Grund og med Videnskabens Vaaben. Dette var allerede blevet
udtalt fra første Færd af. Da Strauß's „Leben Jesu“ var udkommet, forlangte
det preussiske Underviisningsministerium den berømte Neanders Erklæring om,
hvorvidt det maatte være rigtigt at udstede et Forbud mod denne saa megen
Opsigt vækkende Bog. I denne Erklæring hedder det: „De theologiske Læreres
Kald er det at virke til, at den Bog, der er udgaaet fra en den christelige
Kirkes Aand modstræbende Retning, kommer til at bidrage til at fremme
Videnskaben og derved at befordre ogsaa Kirkens egen Interesse. Men dette
Maal kan kun naaes, naar denne Kamp føres med Videnskabens Vaaben
alene“\footnote{Luther derimod bad ved en lignende Anledning de sachsiske
Fyrster holde paa, „daß alleine mit dem Worte Gottes in diesen Sachen
gehandelt werde, wie den Christen gebühret.“}.

Men saa synes jo Skuespillet fra Middelalderens Slutning at gjentage sig?
Her er »\textit{duæ veritates}« --- og begge ville jo være videnskabelige.

Ved en nærmere Betragtning viser dette sig imidlertid ikke at være saa. Vi
ville her kun foreløbigt spørge: hvad er det Strauß og Feuerbach angribe,
Religionen eller Theologien, --- og hvad vil det sige, naar de erklære at
staae paa Videnskabens Standpunkt? --- Det første Spørgsmaal er snart
besvaret:
% p. 17 foot carries the signature "2" (not text).
% --- p. 18 ---
\opage{18}\setcounter{footnote}{0}%
Strauß har foruden „Leben Jesu“, hvor han opløser den evangeliske Historie i
Myther, jo ogsaa skrevet en „Dogmatik“, hvor de dogmatiske Sætninger en
efter en opløses i Intet; og Feuerbach skjelner vel mellem Religion og
Theologie som mellem den umiddelbare, uvilkaarlige Anskuelse og den bevidste
Reflexion, men han vender sine Vaaben mod begge. Der bekæmpes altsaa ikke
blot en Sandhed, der vil være videnskabelig; det er den religieuse Sandhed,
Livssandheden selv der her bekæmpes. Og hvad det andet Punkt angaaer, da
fastholder den negative Kritik de faste og urokkelige Love i Aandens og
Naturens Verden, saaledes som de godtgjøre sig for den menneskelige
Bevidsthed ved Videnskaben. Hvad den rationelle Videnskab ikke kan
controllere og begrunde, det er for dette Standpunkt et \textit{non-ens}.

Saavist som Religion er Et, Theologie et Andet, saavist maa altsaa den
Strauß-Feuerbachske Kritik opfattes under et dobbelt Synspunkt. Overfor
Theologien bliver det en Strid mellem Viden og Viden; overfor Religionen
bliver det en Strid mellem Viden og Tro.

„Striden med Strauß forvandler sig øieblikkelig til en Strid mellem
Religionen og Videnskaben. Det er Kritikernes Fortjeneste at have givet
Sagen et saa stærkt Stød, at den just derved er dreven hen til dette
afgjørende Vendepunkt. Evangelietroens Strid med den straußiske Kritik er en
Competencestrid, en Souverainitetsstrid, en Strid om den øverste Magt, den
høieste Myndighed i Aandernes Verden“\footnote{„Evangelietroen og
Theologien“. Tolv Forelæsninger af R. Nielsen. 1850. \textit{pag.} 85.}.

Det sees altsaa, hvorledes Sagerne stillede sig. Paa den ene Side stod
Dogmatiken med den Erklæring, at Viden stemmede med Troen, paa den anden
Side Kritiken med den ligesaa kategoriske Erklæring, at Viden stred mod og
følgelig opløste Troen. Det syntes virkelig at være paa høie Tid, at
% --- p. 19 ---
\opage{19}\setcounter{footnote}{0}%
Opmærksomheden henvendtes paa „Competencestriden“ mellem Tro og Viden; thi
desangaaende --- og heri vare da Dogmatiken og Kritiken\footnote{Jfr.
Martensens Dogmatik § 17 Anm.} lige dogmatiske --- fik man hos ingen af de
stridende Parter den fornødne Oplysning.

Paa dette Punkt vende vi os nu til de Kierkegaardske Skrifter. Om
nogetsteds, saa er i disse Skrifter Alt opbudt for at fastsætte og hævde
Troesbegrebet i dets Eiendommelighed.

Vi skulle her kun betragte Kierkegaard fra det angivne Synspunkt. Hans store
Betydning gjælder naturligviis for hele Tidsalderen, og hans Optræden blev
da ogsaa bestemt ved Datidens aandelige Tilstand i det Hele. Han er et
sjældent Exempel paa, hvorledes en Begavelse og en Genialitet, saaledes som
den kun bliver faa Mennesker tildeel, udelukkende kan tages i Beslag
(ligesom tages til Fange) af det ethisk-religieuse Formaal: selv
existerende i Sandheden at være Andre en Lærer i at existere i Sandheden.
Men det for ham Særegne ligger i hans Methode --- hans Maieutik, for at
minde om Sokrates. Han opfattede sin Samtid som hildet i et uhyre
Sandsebedrag, i dette nemlig, at „vi Alle ere Christne“; et Sandsebedrag
% PRINTED AS IS: „indrekt“ (a reader has pencilled an "i" with a caret above the word in this copy)
kan aldrig hæves direkte --- derfor „indrekt Meddelelse“! Beroer
Sygdommen paa, at man tager Phænomenet for Væsenet, saa maa Lægemidlet
bestaae i at sætte Splid mellem Væsen og Phænomen, for at Blændværket kan
ophøre og Væsenet erkjendes som Væsen.

Man har ofte indvendt mod denne indirekte Meddelelse, at den var for
„kunstig“, for reflecteret: „Christendommen, Religionen overhovedet er jo
netop betinget ved det umiddelbare Forhold til det Guddommelige!“ Hertil
svarer Kierkegaard: man tager Feil, naar man mener, at det er Lægemidlet,
der indfører Reflexionen; nei, Reflexionen --- det er netop Ulykken --- er
der iforveien, men kun som Halvreflexion,
% p. 19 foot carries the signature "2*" (not text).
% --- p. 20 ---
\opage{20}\setcounter{footnote}{0}%
falsk Reflexion. „Situationen er Christenheden, hvilket er en Reflexionens
Bestemmelse. I Christenheden --- at blive Christen er enten at vorde hvad
man er (Reflexionens Inderlighed, eller Inderliggjørelsens Reflexion), eller
det er først at rives ud af en Indbildning, hvilket igjen er en Reflexionens
Bestemmelse“\footnote{„Synspunktet for min Forfattervirksomhed“,
\textit{pag.} 33.}.

Herpaa beroer da Forklaringen af „Tvetydigheden eller Dupliciteten i hele
Forfatterskabet, om Forfatteren er en æsthetisk eller en religieus
Forfatter“. Det er nu i denne Sammenhæng hverken de æsthetiske eller de
religieuse Skrifter, vi især ville henvise til, men hvad Kierkegaard selv
har betegnet som „Vendepunktet i den hele Forfatter-Virksomhed“:
„Uvidenskabelig Efterskrift“, af Johannes Climacus.

Men hvad vi vel først maa afgjøre, er, om nu ogsaa Johannes Climacus
virkelig mener, hvad han siger. „Han er jo Humorist, spiller og leger med
Alt paa den letfærdigste Maade!“ Lad os da høre, hvorledes han selv
opfattede sin Opgave, da han besluttede sig til at blive Forfatter. „Af
Kjærlighed til Menneskene, af Fortvivlelse over min forlegne Stilling ved
Intet at have udrettet og ved Intet at kunne gjøre lettere end det er gjort,
af sand Interesse for dem, som gjøre Alt let, fattede jeg da det som min
Opgave: overalt at gjøre Vanskeligheder“. „Naar jeg nu turde høirøstet
forkynde, at jeg var kommen til Verden for og kaldet til at modarbeide
Spekulationen, at dette var min dømmende Gjerning, medens min prophetiske
var, at varsle om en mageløs Fremtid, hvorfor Menneskene, paa Grund af, at
jeg var høirøstet og kaldet, med Sikkerhed kunde forlade sig paa hvad jeg
sagde --- saa vilde der vel være Mangen, der i Mangel af at ansee det Hele
for en phantastisk Reminiscents i en Fjantets Hoved, vilde ansee det for
Alvor“\footnote{Hvorledes han saa fik „en nærmere Forstaaelse af sit eget
Indfald, at han maatte see at gjøre Noget vanskeligt“, maa læses hos ham
selv (p. 174 \textit{ff.}).}. Alvor er ikke
% batch 2 ends here: p. 20 ends on the complete word "ikke"; p. 21 continues the sentence.
% --- p. 21 ---
% Batch 3: printed pp. 21--29 (PDF 25--33). p. 21 begins mid-sentence.
\opage{21}\setcounter{footnote}{0}%
saaledes til at tage og føle paa. Alvor er det væsentlige Forhold til det Væsentlige, men ikke ligefrem commensurabel med sit Udtryk. Tvertimod: Alvor er Inderlighed, og „Modsætningens Form er netop Inderlighedens Kraftmaaler“.

Spekulationens Misviisning laa for Johannes Climacus i hele Tidsalderens Retning: „at man overhovedet ved den megen Viden havde glemt hvad det er at existere og hvad Inderlighed har at betyde“. Derfor vilde han „gjøre Vanskeligheder“ --- ved at lægge Eftertryk paa Existentsen. Naar der af det existerende Menneske spørges: „Hvad er Sandhed?“ maa der menes den væsentlige Sandhed, Livs-Sandheden, „den Sandhed, der væsentlig\footnote{%
% PRINTED AS IS: "væsenlig" (first occurrence; t missing), beside "væsentlig" in the next line.
„Kun den ethiske og ethisk-religieuse Erkjenden er derfor væsenlig Erkjenden. Men al ethisk og al ethisk-religieus Erkjenden er væsentlig sig forholdende til, at den Erkjendende existerer.“ Uvidenskabelig Efterskrift p. 145.} forholder sig til Existents“. Sandhed er i Almindelighed Overeensstemmelsen, Eenheden af Subjectivitet og Objectivitet; den indeholder altsaa allerede reent abstract seet en Fordobling. Men denne Fordobling kommer med forstærket Eftertryk igjen, naar det ikke er \textit{in abstracto} men existentielt, at der spørges om Sandheden. Her bliver Tvesidigheden til Disjunction. Erkjendelsens to Veie gaae i modsat Retning: den ene vender sig mod det Objective, og fører --- netop for ret at faae det Objective saa objectivt frem som muligt --- bort fra det Subjective; den Anden fører til Subjectiviteten, til Sandhedens Tilegnelse i Inderlighed.

Den subjective, existerende Tænkers Thesis bliver da: Subjectiviteten er Sandheden. Der er en Sphæreforskjel mellem den videnskabelige og den existentielle Sandhed. Her er ingen Mediation mulig; thi ved Mediationens Subject-Object kommer man tilbage til Abstractionen, medens Spørgsmaalet jo netop var, „hvorledes et existerende Subject \textit{in concreto} forholder sig til Sandheden.“

% --- p. 22 ---
\opage{22}\setcounter{footnote}{0}%
„Men“ --- siger man --- „det er dog høist eensidigt at sige, at Subjectiviteten er Sandheden; Subjectiviteten maa jo have optaget det Objective for at være Sandheden.“ Denne Indvending beroer deels paa en „ufuldkommen og udialectisk Distinction“ mellem Subjectivt og Objectivt, deels har den ikke forstaaet den anførte Thesis ret. Lad os høre Johannes Climacus. „Medens den objective Tænkning er ligegyldig ved det tænkende Subject og dets Existents, er den subjective Tænker som existerende væsentligen interesseret i sin egen Tænkning, er existerende i den. Derfor har hans Tænkning en anden Art af Reflexion, nemlig Inderlighedens, Besiddelsens, hvorved den tilhører Subjektet og ingen Anden . . . Inderlighedens Reflexion er den subjective Tænkers Dobbelt-Reflexion“\footnote{Uvidenskabelig Efterskrift. p. 50---51.}. Der er altsaa ikke blot et Reflexionsforhold mellem det Objektive og Tænkningen, men ogsaa mellem det tænkte Objective og Tænkerens Existents. Men denne sidste Reflexion er ikke logisk, men practisk, er et \textit{salto mortale}, for at minde om Jakobi, en „dialectisk Gjerning og Daad“\footnote{Stilling: „Om den indbildte Forsoning af Tro og --- Viden“. 1850. p. 30---32.}.

Hvad der hidtil er udviklet, angiver Vilkaarene for enhver Existerendes Forhold til Sandheden: det er „den socratiske Viisdom, hvis udødelige Fortjeneste er netop at have agtet den væsentlige Betydning af, at den Erkjendende er existerende, hvorfor han i sin Uvidenhed i høieste Forstand indenfor Hedenskabet var i Sandheden“\footnote{Uvidenskabelig Efterskrift. p. 150.}.

Skal der da gives et fra det sokratiske væsentligt forskjelligt Standpunkt, maa denne Forskjel beroe paa den nærmere Bestemmelse af Sandheden.

Den Sætning: „Subjectiviteten er Sandheden“ indeholder en Modsætning til
% PRINTED AS IS: "Objectiviten" (for Objectiviteten?).
Objectiviten. Denne Modsætning maa da ogsaa finde sit Udtryk i Sandhedens Bestemmelse. Medens
% --- p. 23 ---
\opage{23}\setcounter{footnote}{0}%
den objective Tænkning søger Visheden i de objective Grunde, sees Sandheden fra det subjective, Existentsens Standpunkt at være objectiv Uvished. De Grunde, hvorom her kan være Tale, ere blot subjective, existentielle.

„Naar Subjectiviteten, Inderligheden er Sandheden, saa er Sandheden objectivt bestemmet Paradoxet; og det at objectivt Sandheden er Paradoxet, viser netop, at Subjectiviteten er Sandheden“\footnote{Anf. Skr. p. 151.}.

Paradoxet er tilstede, hvor det Evige og Ubetingede træder overfor det Timelige og Betingede. Ogsaa for Sokrates var der et Paradox. „Dog er den evige væsentlige Sandhed selv ingenlunde Paradoxet, men er det ved at forholde sig til en Existerende“. Thi at det Uendelige forholder sig til det Endelige, er vel paa Existentsens Standpunkt en Modsigelse, som kun overvindes ved et \textit{salto mortale}, men paa Videnskabens, Begribningens Standpunkt medieres denne Modsigelse\footnote{„Ikke ved et Spring men ved en til Functionernes Eiendommelighed svarende Overgang maa den begribende Subjectivitet i hvert Tilfælde bevæge sig fra det Endelige til det Uendelige“. Grundideernes Logik p. 230.}.

% Greek accented as printed: grave on the alpha, circumflex on the eta.
Men hvis nu den evige Sandhed og det at existere i Timelighed sættes sammen i Sandheden selv, saa bliver Sandheden selv paradox, --- og saa haves Paradoxet κὰτ᾽ ἐξοχῆν, uden at Mediation er mulig.

Her staae vi da ved Aabenbaringen. Den evige, personlige Gud taler til Mennesket i Tiden: „Og Gud fristede Abraham, og sagde: Abraham, Abraham hvor er Du? men Abraham svarede: her er jeg“ (1 Moseb. 22, 1). Gud selv fødes og lever og døer som Menneske i Tiden: „Ordet var hos Gud, og Ordet var Gud. Og Ordet blev Kjød og boede iblandt os“ (Joh. Ev. 1, 1. 14). Her hjælper al historisk og spekulativ Viden Intet; det Høieste, der ad den objektive Vei
% --- p. 24 ---
\opage{24}\setcounter{footnote}{0}%
% PRINTED AS IS: „den objektive Vished (paa det sokratiske Standpunkt)“ (confirmed at the image, 23 Sep 2026). The quoted Postscript reads „den objektive Uvished“, as does p. 23 here; the translation renders "uncertainty".
kan naaes, er en Approximation, altsaa ingen qvalitativ Afgjørelse, hvilken kun kommer ved Troen. „Istedetfor den objektive Vished (paa det sokratiske Standpunkt), er her Visheden om, at det objektivt seet er det Absurde, og dette Absurde fastholdt i Inderlighedens Lidenskab\footnote{„Lidenskaben er netop det Spændende i Modsigelsen; naar denne tages bort, er Modsigelsen et Plaisanterie, et Bonmot“. Uvid. Efterskr. p. 294.} er Troen.“

Johannes Climacus har da her opstillet et Troesprincip, der er uforeneligt med Viden. Og det er ikke alene fra først af, naar Subjectiviteten vælger Troen, at denne frastøder og frastødes af Viden; hvad der principielt gjælder for Udviklingen i dens Begyndelse, maa ogsaa gjælde under dens videre Gang.

Dette sit Princip stiller Johannes Climacus nu mod de to forhen angivne Retninger. Mod Strauß og Feuerbach hævder han, at Videnskaben mangler Competence til at dømme paa Existentsens Omraade; mod den videnskabelige Religionslære, at Troens blanke Skjold ikke maa besmittes med nogen Plet af objectiv Viden. At hans Polemik hovedsagelig rettede sig mod den sidste, laa vel for en Deel i, at den var den raadende hertillands, hvor den negative Kritik ingen originale Repræsentanter havde; men fornemmelig var det, fordi den efter hans Mening paa en uklar og tvetydig Maade sammenblandede den moderne Kultur og Christendommen --- heri i Overeensstemmelse med Tidens Retning i det Hele.

Vi have da her Synet af tre Modstandere, som gjensidig bekrige hverandre. Kunde der maaskee være Tale om et Forbund mellem de To af dem mod den Tredie, vilde den dog aabenbart kun blive kortvarig.

Opmærksomheden retter sig naturligt mod den af Parterne, der maaskee kunde synes at forene de tvende andre i en „høiere Eenhed“: Dogmatiken, som jo \emph{baade} hævdede Tro \emph{og} Viden. Det er klart, at der nu var Forudsætninger tilstede for en
% --- p. 25 ---
\opage{25}\setcounter{footnote}{0}%
rigere og dybere Udvikling af denne Videnskab: paa den ene Side var jo det subjective Troesprincip, paa den anden Side det objective Vidensprincip gjennemført saa skarpt og conseqvent, som tænkes kan. Her var en Seir at vinde for den, der kunde magte disse Modsætninger.

Da udkom i Aaret 1849 Martensens „christelige Dogmatik“: „en videnskabelig Fremstilling og Begrundelse af de christelige Troeslærdomme i deres indre Sammenhæng“ (§ 1).

For den, der mindedes, hvilke Fordringer Forfatteren selv tidligere havde stillet til en videnskabelig Dogmatik, --- hvorledes han, i Anledning af det af Mynster opstillede „Enten-Eller“ mellem Supranaturalisme og Rationalisme, havde hævdet Mediationen og erklæret\footnote{Se „Tidsskrift for Litteratur og Kritik“ 1839, især p. 464, 469.}, at en Tilbagevendelse til den christelige Aabenbaringssandhed uden som en „nødvendig, i vor Tids Philosophie selv liggende Vending“, fremgaaet af en Gjennemførelse af Hegels dialectiske Methode, var uvidenskabelig (eller dog „ikke noget videnskabeligt Fremskridt“), --- for den var der al Opfordring til at undersøge den nye „christelige Dogmatik“ med særligt Hensyn til dens videnskabelige Princip. Var hiint Program fastholdt, og var der endog blot prøvet paa dets Udførelse, da vilde det begynde at see alvorligt ud for den negative Kritik: det var jo lovet, at i den „immanente Erkjendelse af Christologien“ skulde ogsaa den Straußiske Christologie være optaget som en „relativ Opfattelsesmaade“, en „ophævet Idealitet“.

For den, der havde seet, med hvor megen Standhaftighed og Vedholdenhed en heel Række af pseudonyme Forfattere, hvis Formaal var at hævde den existentielle Sandhed og Troen som dens høieste, paradoxe Form, havde udslynget den ene Udfordring efter den anden til „Systemet“, som vel ikke endnu var „færdigt“, men dog forventedes snart at skulle blive det, --- for den var der al Opfordring til at prøve, hvorledes
% --- p. 26 ---
\opage{26}\setcounter{footnote}{0}%
den nye „christelige Dogmatik“ vidste at forlige Paradoxet med Videnskaben, den existentielle med den videnskabelige Erkjendelse.

„Det Spørgsmaal“, hedder det hos den første af de Forfattere, der optraadte mod den Martensenske Dogmatik\footnote{„Mag. S. Kierkegaards „Johannes Climacus“ og \textit{Dr.} H. Martensens „Christelige Dogmatik“. En undersøgende Anmeldelse af R. Nielsen“. 1849.}, „hvorom Theologien og Philosophien for nærværende Tid maae siges at bevæge sig, som om det sande og egentlige Principspørgsmaal, kan kortelig fremsættes i de Ord: Vil Christendommens Sandhed efter sin Natur være Gjenstand for objectiv Viden?.. Som Principspørgsmaal kan det paa ingen Maade holdes hen i undvigende Ubestemthed, men maa besvares ubetinget, med Ja eller Nei.“ Forfatteren sammenstiller derpaa de to forskjellige Besvarelser, som Pseudonymen Johannes Climacus og Forfatteren af den „christelige Dogmatik“ give: den ene med „et ubetinget afgjørende Nei“, den anden --- „skjøndt dog med nogen Usikkerhed“ --- med Ja. Mod denne sidste Besvarelse rettes derpaa Kritiken, for at vise, at „den „christelige Dogmatik“ udialectisk har omgaaet det christelige Problem“, medens Pseudonymen skjærper dette. Der vises først, at det ikke er lykkedes Martensen at adskille \emph{Dogmatik} og \emph{Religionsphilosophie}, „theologisk“ og „philosophisk Spekulation“, men at han tvertimod modsiger sig selv (smlgn. Dogmatikens § 2 med § 8 Anm.). For nu nærmere at oplyse denne Uklarhed rettes Øiet paa „Spekulationens Problem“; der spørges, om den „christelige Dogmatik“ anerkjender „\emph{de Fordringer, den objective Tænkning paalægger den Spekulerende}“, hvilke hovedsagelig ere: „et directe Forhold mellem den theoretiske Vished og den spekulative Indsigt“, og, hvad Methoden angaaer, „det dialectiske Forhold mellem Forudsætningen og Abstractionen“. Disse Fordringer vilde kræve, at alle Dogmer forvandledes til Problemer\footnote{%
% RUN-OVER FOOTNOTE: this note begins on p. 26 and continues on p. 27, where
% the note area opens without a mark; the carry-over point is marked below.
Cfr. „Tidsskrift for Litteratur og Kritik“ 1839, p. 464, hvor Dogmatikens Forf. erklærede, at „Videnskaben ikke tilfulde kunde
% [the note is carried over to p. 27 at this point]
forsones med noget Dogma, førend det saaes i dets dialectiske Nødvendighed i Systemet“.}. Men derpaa
% --- p. 27 ---
\opage{27}\setcounter{footnote}{0}%
kan den „christelige Dogmatik“ ikke indlade sig: Christendommens absolute Sandhed er for den urokkelig: den \emph{troer} paa den. Derfor henvender da Kritikeren for det Tredie sin Opmærksomhed paa „Troens Problem“, og da navnlig paa \emph{Forholdet mellem Troen og dens Gjenstand}. Dette Forhold er ifølge Dogmatiken et dobbelt: et existentielt og et theoretisk Forhold. Vel --- siger Kritikeren --- opstil da den objective \emph{Troesgjenstand}, Daaben og Nadveren f. Ex., og lad os --- i Troen --- faae en videnskabelig-objectiv Afgjørelse af, hvad der er den sande Daabs- eller Nadverlære! Men se: den videnskabelige Theorie afhænger af, om det er en \textit{theologia regenitorum} eller \textit{irregenitorum} (Rationalister), og hvad den første angaaer, igjen af, om den Gjenfødte er Lutheraner (og her igjen: Orthodox eller Pietist, „Skrifttheolog“ eller Grundtvigianer) eller Reformeert o. s. v. Det sees da: hvorledes Troens objective Gjenstand er, beroer paa Troen, paa det Subjective. Og gaaer man saa ud fra \emph{Troen}: hvorledes kan da „en Erkjendelse, der er saa subjectiv, saa inderlig, saa absolut personlig, som en sand Troeserkjendelse nødvendigviis maa være, omsættes i Videnskabens objective, upersonlige Form og endda være Troeserkjendelse?“ Herom gav Dogmatiken ingen Oplysning. Eftersom altsaa den „christelige Dogmatik“ \emph{hverken} har opstillet og hævdet et særegent, af Philosophien uafhængigt, theologisk Erkjendelsesprincip, \emph{eller} har vedkjendt sig og gjennemført den objective Tænknings Princip, \emph{eller} har resolut og uden Halvhed stillet sig under Troens Banner eller dog fremstillet Troesprincipet reent og skjært, --- saa ender Kritiken med den Thesis: „Den christelige Dogmatik“ er principløs\footnote{Hvilken Kritikerens Dom var, naar han saae bort fra Principet, og det blot „kom an paa at bedømme dette omfangsrige Arbeide med udelukkende Hensyn til Forfatterens Talent, Studium og Fremstillingsgave“: se den „undersøgende Anmeldelse“ p. 42---44.}.

% --- p. 28 ---
\opage{28}\setcounter{footnote}{0}%
% PRINTED AS IS: "eu" (turned n) for "en" in "med eu Kritik".
% Quotation marks: »…« (guillemets) round the Latin tags on pp. 28--29, beside the usual „…“.
Den anden Forfatter, der optraadte med eu Kritik af den Martensenske Dogmatik, opstiller i sin Afhandling\footnote{„Om den indbildte Forsoning af Tro og --- Viden, med særligt Hensyn til Prof. Martensens „christelige Dogmatik“. Af Mag. P. M. Stilling. 1850.} tvende forskjellige Forhold mellem Viden og Tro. Det første er: »\emph{\textit{qværo intelligere ut credam}}«. Her er da Troen ikke Forudsætning. Den efter Erkjendelse higende Aand gjennemvandrer med sin intellectuelle Kjærlighed Objectivitetens forskjellige Sphærer, og aner „i Erkjendelsesprocessens Iver og i den theoretisk-reflecterende Aands første Arbeidsdage“ ikke, at den i den religieuse Sphære er anderledes situeret end i de andre Sphærer. „Vel gjælder ogsaa i den religieuse Sphære den Hegelske Bestemmelse af Sandheden som Proces, Synthese af det Subjective og Objective . . . . Men --- vel at mærke og ikke at forglemme --- \emph{saaledes}, at det Objective i denne Sphære kun subjectiveres, tilegnes og inderliggjøres i Tro og Haab, men ingenlunde i spekulativ-theoretisk Proces“. Snart kommer det derfor for en Dag, at \textit{pro} og \textit{contra} her staae mod hinanden, --- at der er „tvende Veie og ingen Vished forud“. Men førend Subjectiviteten beslutter sig til at gribe den eneste Vished, her kan være Tale om: Troens og Haabets, vil der maaskee finde alvorlige theoretiske og practiske Kampe Sted. „Og mindst af Alt kunde en Reflecterende af „hine tappre Rækker“ falde paa at lade sin Tankekraft engagere som Dogmernes Jabroder; et Engagement, der vel stundom tilbydes den af dogmatiserende Tænkere“. --- Den Sætning, at Subjectiveringsprocessen i den religieuse Sphære er qvalitativt forskjellig fra den theoretiske „forglemmes i meer eller mindre Grad hos Hegel, Marheineke, kort hos alle spekulative Theologer. Det nytter slet ikke, at Prof. Martensen siger: „En \emph{Troende} kan tilegne sig det Christelige i spekulativ-theoretisk Proces. Det er grundfalsk“. For at vise dette gaaer Forf.
% --- p. 29 ---
\opage{29}\setcounter{footnote}{0}%
over til det andet Standpunkt: »\emph{\textit{credo ut intelligam}}«. Dette er „Reflexionen \textit{post fidem} (ɔ: efterat jeg har taget Troespartiet)“. Men --- siger Kritikeren --- „det gaaer slet ikke bedre med Reflexionen \textit{post fidem}, end \textit{ante fidem} (ɔ: førend jeg tog Troespartiet og endnu stak lige midt i Reflexionsmodsætningerne). Reflexionsmodsætningen (Vantro) blev i første Troesskridt overvunden, ikke ved Reflexionen, men ved og i troende Leven, der dicterede Lydighed . . Og ethvert følgende Skridt (2. 3. 4 . . .) er af samme qvalitative Natur, som \emph{første} Skridt, der idelig repeteres under den Troendes og Haabendes videre Gang. Og vel maa der tales, og tales ogsaa i Skriften, om Fremgang og Fremvæxt i Tro og Haab, nemlig i den subjective, personlige Visheds uendelige Styrke. Men der tales sikkerlig ikke om, i Fremvæxt at „voxe ud over“ eller i „Fremgang“ at gaae ud over Troes- og Haabsvishedens qvalitative Grændser til Bevisets \emph{objective} Vished“\footnote{Anf. Skr. p. 37 f.}. Da imidlertid den „christelige Dogmatik“ netop havde opstillet et saadant »\textit{credo ut intelligam}« som sit Princip, spørger Kritikeren: „Er det virkelig lykkedes Prof. Martensen at hæve Tankekorset og erholde Reflexionskraftens oprigtige Understøttelse og Sanction af Dogmet?“ Og det Resultat, hvortil Kritikeren saa kommer efter en Prøvelse af Dogmatikens Behandling af de vigtigste Dogmer, er „at »\textit{contra}« slet ikke er beseiret, men kun overhørt eller neddysset, at Reflexionsvidnesbyrdet kun partisk er afhørt, men aldrig faaer Lov til at tale ud“, „at vor Dogmatikers »\textit{credo}« i Gjerningen kun er et scholastisk hemmende \textit{credo}, der lægger Reflexionsfriheden under Censur og ikke tilsteder den andre Yttringer, end saadanne, der ere Dogmatikerens Øre velbehagelige“ (p. 66). ---
% PRINTED AS IS: »contra“ -- opened with a guillemet, closed with a high double comma.
Og Kritikeren finder noget Teleologisk i, at »\textit{contra}“ ikke videnskabeligt kan skaffes tilside. Gaves der nemlig et „Beviis \textit{post fidem}“, saa vilde dette have en uheldig Tilbagevirkning
% p. 29 ends here, mid-sentence, after the whole word "Tilbagevirkning".
% ==== Batch 4: printed pp. 30--38 (PDF 34--42). Folios read at the head of every leaf.
% Conventions in this fragment: quotation marks as printed, „ … “ (low-high).
% Fraktur ſs/ß in Danish words (Interesse, disse) → ss; ß kept in German words
% and in „Strauß“ (corpus convention). „p.“ before page numbers is antiqua but
% left unmarked (corpus convention); Dr., Latin phrases and pro/contra, set in
% antiqua, → \textit{}. Printed number ranges with a rule → "--".
% --- p. 30 ---
\opage{30}\setcounter{footnote}{0}%
paa Troes- og Haabsvisheden. „Det er nemlig psychologisk
umuligt, ja kun en pompøs Talemaade, at jeg paa eengang
har Troes- og Haabsvisheden „uafhængigt af Spekulationen“,
„ad en heel anden Vei end Spekulationen“, og dog tillige i
samme Aandedrag just ad Spekulationens Veie mere og mere
(Grændsen flyder . . .) faaer Troes- og Haabsvisheden
confirmeret“ (p. 64 Anm.).

Det kan ikke negtes, at Kritiken havde angivet klart og
bestemt, hvad det var det kom an paa. Hvad Indsigelserne
rettede sig imod, var ikke dette eller hiint Enkelte, men det
var selve det dogmatiske Princip. Den Fordring maatte derfor
nødvendig gjøres til Antikritiken, om der kom en saadan, at
den ikke --- dogmatisk --- ansaae det dogmatiske Princip for
at staae urokkeligt fast, men indlod sig paa en nærmere
Begrundelse deraf, samt udviklede klart og utvetydigt dets
Forhold til den „rene“ Viden paa den ene Side, og til Troen
paa den anden Side.

% footnote mark 1) on p. 30
Men det Svar\footnote{„Dogmatiske Oplysninger“. Et Leilighedsskrift af \textit{Dr.} H. Martensen. 1850.}, Dogmatikens Forfatter gav sine
Kritikere, var nu netop et reent dogmatisk. Vel er det „en
gammel Erfaring, at Dogmatiken altid har Kritiken i sit
Følge“, og „især i vor Tid, der er langt mere kritisk end
dogmatisk“; men hvad Kritiken over nærværende Dogmatik
angaaer, da mener Forf., at den „endnu ikke kan siges at
være begyndt“ (p. 7). Dette kunde synes underligt for den,
der har læst de to ovenfor omtalte kritiske Afhandlinger; men
ved nærmere Undersøgelse opdager man, at Kritiken afvises,
fordi den angriber selve Principet: naar den gjør det, saa er
den fra det dogmatiske Standpunkt ikke „rigtig“ Kritik, ja
mærkes egentlig slet ikke af Dogmatiken. Saaledes stiller
denne sig f. Ex. overfor Strauß, og overfor
% PRINTED AS IS: „euhver“ for „enhver“ (turned n; checked at 800 dpi against „nu“ on the same page).
euhver „gjennemgribende Tvivl om Christendommens Grundlag“; hvor „be%
% --- p. 31 ---
\opage{31}\setcounter{footnote}{0}%
tydningsfuld“ en saadan Tvivl ellers i andre Henseender kan
være: den vedkommer slet ikke Dogmatiken! Og hvorfor ikke?
Fordi „den dogmatiske Interesse“ slet ikke existerer for en
saadan Tvivler!\footnote{Se Martensens Dogmatik § 31 Anm.}

Her er da sat en forsvarlig Bom for enhver Kritik.
En egentlig Forhandling mellem Dogmatiken og Kritiken
kunde der altsaa ikke være Tale om. Men forøvrigt har den
„dogmatiske Oplysning“ en forskjellig Characteer.

Overfor Mag. Stilling f. Ex. viste det „Dogmatiske“ sig
i, at „Hr. Magisterens Skriveri“ strax blev „lagt hen for ikke
mere at optages“ (p. 8). Med Rette yttrede Mag. Stilling
herom i en senere Afhandling: „Ved hvilke uovervindelige og
høist ærefulde Vaaben en spekulativ Dogmatiker formaaer at
haandhæve sit Standpunkt mod udførligen begrundede Angreb
fra fordums Venners Side, derpaa afgiver Prof. Martensens
Banstraale over min Kritik, et meget belærende Exempel.“

Den vægtigste Kritik over Dogmatiken var allerede skrevet,
før denne udkom: af Pseudonymerne. Desuagtet var der i
selve Dogmatiken slet intet Hensyn taget til disse Forfattere,
med Undtagelse af, at det i Forordet bebreides „Nogle iblandt
os“, at de ville være „absolute i Troen“\footnote{Som man let seer, et tvetydigt Udtryk: thi Et er det at hævde
Troesprincipet absolut, uden at slaae af, et Andet er det at erklære
sig selv for at staae paa Troens absolute Høide. --- Iøvrigt var
der i Forordet ogsaa hentydet til „nogle Enkelte, som uden at føle
Drift til sammenhængende Tænkning, kunne tilfredsstille sig selv ved
at tænke i Strøtanker og Aphorismer, Indfald og Glimt“. Saaledes
lød Dommen 1849. Mon den ogsaa vilde lyde saaledes 1866?}. Nu havde Prof.
Nielsen i sin Kritik fremdraget Johannes Climacus igjen, for
at det Forsømte mulig kunde gjenoprettes. De „dogmatiske
Oplysninger“ maatte følgelig oplyse Dogmatikens Forhold til
% PRINTED AS IS: „Pseydonymerne“ (y for u; the same word is spelt „Pseudonymerne“ six lines above).
Pseydonymerne. Uagtet nu disse Forfattere erklæres for „i
nærværende Sammenhæng aldeles uvedkommende“, uagtet Forf.s
% --- p. 32 ---
\opage{32}\setcounter{footnote}{0}%
„Kjendskab til denne vidtløftige Litteratur kun er saare ringe
og fragmentarisk“, erklærer han dog, at han af disse Skrifter
har en heel anden Opfattelse end Recensenten: den nemlig: at
de ikke vilde „begrunde en Reform af Dogmatiken“ men „paa
sokratisk Viis sætte en dybere Skepticisme i Bevægelse, for
derved at opvække Læseren til at søge det Problem, som er
høiere og Mere end alle theologiske og philosophiske Skoleproblemer,
nemlig det personlige Livsproblem“ (p. 12). Og
man behøver ikke at have læst meget i de pseudonyme Skrifter
for at sande disse Ord. Men for det Første var Forf. af
Dogmatiken heri slet ikke uenig med sin Recensent; thi denne
havde jo netop villet forebygge, at hans „Opfattelse af Bogen
(ɔ: „Uvidenskabelig Efterskrifts“) videnskabelige Betydning forvexles
med Forfatterens derfra maaskee noget afvigende\footnote{Hermed kan jevnføres det for Kierkegaard saa characteristiske Forhold,
hvori han stillede sig til denne Polemik. Se „Kierkegaards
Bladartikler“ p. 252 Anm.}
% PRINTED AS IS: no closing quotation mark after „undersøgende Anmeldelse“.
Hovedtanke“ (Prof. Nielsens „undersøgende Anmeldelse p. 7).
Og dernæst --- maa ikke „Troesvidenskaben“ falde, saafremt
Johannes Climacus's Troesbegreb er rigtigt? Med mindre
man da anseer den „dybere Skepticisme“, han vil sætte i
Bevægelse, for saa dyb, at den ingen virkelig Betydning
faaer. --- For øvrigt ville de „dogmatiske Oplysninger“ ikke
have Noget med Joh. Climacus selv at gjøre; han afvises som
uvedkommende; men kun med Prof. Nielsens Fremstilling af
hans --- „Sætninger“. Det erklæres da, at „det Nye i dem
er ikke sandt, det Sande ikke nyt“; der omtales derpaa endeel
gamle Sandheder --- uden at der undersøges, om det Nye
maaskee ikke netop laa i Maaden, hvorpaa de fremsattes\footnote{„Objectivt accentueres: \emph{hvad} der siges; subjectivt: \emph{hvorledes} det
siges“. Uvidensk. Eftersk. p. 149.};
indtil Forf. standser ved to „nye Sætninger“ for at bevise,
at de ikke ere Sandheder. Den første af disse: „Sandheden,
Inderligheden og Subjectiviteten“, var Forf. ny „i den For%
% --- p. 33 ---
\opage{33}\setcounter{footnote}{0}%
stand, at han ikke havde fundet den i nogen ham bekjendt
Dogmatik“ (p. 16). Dette var endda ikke saa underligt, da
Sætningen (idetmindste i sin Bitanke, sin Conseqvents) er
polemisk mod den videnskabelige Dogmatik. Den anden Thesis,
som var Forf. ny, var: „for den existerende Subjectivitet er
den høieste Sandhed Paradoxet“. I denne Sætning tænkte
Forf. med Rette, at „Problemet maatte være at finde“; men
uheldigviis glemte han, at man for at erkjende, at Sandheden,
den høieste Sandhed, er Paradoxet, nødvendigviis maa gaae
ud fra den \emph{existerende} Subjectivitets Standpunkt: hvad
Under da, at hans objectiv-dogmatiske Undersøgelse kommer til
det Resultat at „paradox“ kun er en rhetorisk Hyperbol, og
Paradoxet Noget, der blev opfundet af Slangen i Paradiis
(p. 23)!

Spørger man, hvad de „dogmatiske Oplysninger“ lære
om Forholdet mellem den existentielle og videnskabelige Erkjendelse,
da er der især to Punkter, vi kunne henlede Opmærksomheden
paa. For at vise, at Troen i sig indeholder
en Spire til Videnskab, indskærper Forf., at Troen ikke blot
er „Forløsningsbevidsthed“ men „Aabenbaringsbevidsthed“ (p.
30--40). I ikke at agte paa dette skal Grundfeilen i Johannes
Climacus's, af Prof. Nielsen fremsatte, Troesbegreb bestaae.
Men hvad i al Verden hjælper denne Oplysning? Bliver
Troen mere videnskabelig, fordi den er Aabenbaringsbevidsthed?
Eller er ikke Aabenbaringen netop kun sand for den \emph{Troende},
og kun \emph{forsaavidt, som} han er troende? Er det en \emph{videnskabelig}
Stræben „i Forløsningen at søge Aabenbaringen“,
„ad Andagtens og den troende Betragtnings Vei at søge at
erkjende den Gud, som gjennem sit Forløsningsværk afslører
sit Væsen for vort Blik“? Men lad os alligevel for et Øieblik
antage, at det lykkes Dogmatiken at faae begyndt paa sin
Opgave: „videnskabeligt at udvikle Troens contemplative Moment“,
--- af hvad Beskaffenhed bliver da den videnskabelige
Erkjendelse, der deraf resulterer? Her hævder Forf. --- og
% (signature mark „3“ at the foot of p. 33, not transcribed)
% --- p. 34 ---
\opage{34}\setcounter{footnote}{0}%
dette er det andet Punkt, vi ville fremdrage --- at „den dogmatiske
Grandskning og Ransagen af Aabenbaringens Sandheder“
ikke er „en Afart eller Forvanskning af den philosophiske
Spekulation“, men har „sit eget, af Philosophien uafhængige
Princip, og derfor ogsaa bør dømmes efter sin egen Maalestok“
(p. 66--68); thi „Forløsningen af den faldne Menneskenatur“
maa ogsaa omfatte Tankens Forløsning, saasandt det
er „Christendommens Hensigt, fuldstændigt at restituere den
sande Humanitet“ (p. 22).

Der gives altsaa et specifisk theologisk-videnskabeligt Erkjendelsesprincip
--- hvilket ikke skal være „Andet, end hvad
der ligger i den ret forstaaede Lære om \textit{testimonium spiritus
sancti}“ (Dogmatiken § 29). Dette er netop et Kjærnepunkt
i den hele Strid. Prof. Nielsen havde allerede rettet sin Opmærksomhed
herpaa i den „undersøgende Anmeldelse“; i den
„Belysning“ af de „dogmatiske Oplysninger“, hvormed han
sluttede Polemiken\footnote{„\textit{Dr.} H. Martensens „Dogmatiske Oplysninger“ belyste af Prof. R.
Nielsen“. 1850.}, hedder det med Hensyn til den omtalte
„theoretiske Betydning af \textit{testimonium spiritus sancti}“: „Der
skal altsaa, ifølge \textit{Dr.} Martensens „Grundanskuelse“ gives
tvende Arter af objektiv Erkjendelse, nemlig: en lavere, verdslig,
almeenvidenskabelig, der endnu ikke kan sige: \textit{Credo ut intelligam},
efterdi den endnu er i sin Syndighed, . . . og en
høiere, kirkelig, dogmatisk-videnskabelig Erkjendelse, der strax
kan sige: \textit{Credo ut intelligam}, efterdi den er forløst ved Guds
Naade, og det saaledes forløst, at den kan producere den ægte
theologiske Spekulation . . . . I den spekulative Prædikestiil,
der prædiker Spekulationen uden at tænke Tanken, kan en
saadan Forsikkring om „Tankens Forløsning“ vistnok tage sig
ret godt baade opbyggeligt og spekulativt ud; men --- „viis
mig din Tro af dine\footnote{Jfr. Grundideernes Logik I, p. 44--53, hvor det i et enkelt Exempel
vises, hvorlidt Gjerning her svarer til Tro.}
% PRINTED AS IS: „37--38“ set with two short hyphens, not the rule used in „66—68“ above.
Gjerninger!““ (p. 37-{}-38).

% --- p. 35 ---
\opage{35}\setcounter{footnote}{0}%
Man maa nemlig spørge: har man en anden Logik og
Ontologie paa det dogmatiske Standpunkt (ɔ: efterat være
bleven gjenfødt) end ellers i Videnskaben? Naar det nu er
klart, at Theologien i denne Henseende er de samme Vilkaar
underlagt som de almeenmenneskelige Videnskaber, hvor skal
saa denne Logik og Ontologie controlleres? Mon ikke i Philosophien,
der „ved sine logiske og ontologiske Undersøgelser
danner et propædeutisk Grundlag for al Videnskab“ (Martensens
Dogmatik § 36)? Men hvorledes kan den „ikke troende“, „ikke
gjenfødte“ Videnskab da controllere den „troende“ og „gjenfødte“?
Og hvad andet Udfald vil denne Controlleren have, end at
der maa anvises den „\emph{troende} Videnskab“ „en egen Plads
\emph{udenfor} Videnskaben“? Og hvis man saa vilde svare herpaa
med at citere 1 Kor. 2, 14--15, saa vilde det paa eengang
være en Bekræftelse paa, at denne Dom var rigtig, og en
stor Misforstaaelse af, hvad „Guds Aand“ betyder.

Et Princip skal vel ikke søge sin Begrundelse udenfor sig
selv, men være et Selvgyldigt. Men et videnskabeligt Princip
maa videnskabeligt kunne gjøre Rede for sig selv. Det nytter
ikke, naar Controllen bliver nærgaaende, da at appellere til
Troen og erklære, at man paa Basis af den har en høiere
Videnskabelighed.

Og hermed slutte vi da Fremstillingen af denne Polemik.
% UNSURE: „litteraire“ — the first e after „litt“ is blotted/damaged in the print; e is the best reading.
Det er ofte blevet sagt, at saadanne litteraire Polemiker ere
ufrugtbare og resultatløse. Og ganske vist kan det være vanskeligt
i endelig Forstand at angive det Resultat, hvortil man
er kommet. Men ved enhver dygtig Polemik opnaaes dog idetmindste,
at de stridende Principer komme til Orde og træde
frem paa Kamppladsen for at vove en Dyst med hinanden. Og
mon det dog saa ikke snart vil blive aabenbart, hvem der,
ideelt seet, er den Overvundne? Naar der er opstillet Indsigelser,
som ramme Sagen i dens inderste Væsen, og man
saa ligefrem afviser dem uden at indlade sig ret med dem, er
det saa ikke et Nederlag?
% (signature mark „3*“ at the foot of p. 35, not transcribed)

% --- p. 36 ---
\opage{36}\setcounter{footnote}{0}%
Og netop paa denne Maade optraadte Dogmatiken, i den
her omhandlede Polemik, mod Kritiken. Det tænker jeg at
have viist ved ovenstaaende Fremstilling. Vil man kalde den
partisk og eensidig, saa svarer jeg: viis mig en anden Maade
at opfatte Sagen paa; viis, at Dogmatiken virkelig har indladt
sig med og overvundet Kritiken! viis, at Dogmatiken har
kunnet møde med god Samvittighed saavel for Videnskabens
som for Troens Domstol, at den har kunnet forene begge
uden at halvere og slaae af paa dem begge!

De tre Kæmper, der tidligere stode paa Kamppladsen
mod hverandre, ere saaledes blevne reducerede til to. Men
denne Tvekamp giver Anledning til adskillige Betænkeligheder.
Paa den ene Side staaer Troen, paa den anden Side --- ja
hvad staaer nu egentlig der? Strauß hævder den moderne Bevidsthed
og dens Verdensanskuelse i Modsætning til „den gamle,
især orientalske Verdens“ Standpunkt, hvorfra de bibelske
Skrifter ere affattede. „Den nyere Tid har en Række af
møisommelige Granskninger, der ere fortsatte gjennem Aarhundreder,
at takke for den Indsigt, at Alt i Verden hænger
sammen ved en Kjede af Aarsager og Virkninger, der ingen
Afbrydelse taaler . . . . Fra dette Standpunkt af kunne de
bibelske Fortællinger umulig vise sig som Historie . . . . Det
Guddommelige kan ikke være skeet saaledes; det saaledes Skete
% (the g of „guddommeligt“ is partly broken in the print; reading not in doubt)
kan ikke have været guddommeligt.“ Hvorledes stiller da her
den moderne, videnskabelige Bevidsthed sig overfor Christendommen?
Ja, den gjør Alvor af, at Religionsforskeren skal
forholde sig til Aabenbaringen som Naturforskeren til Naturen:
nemlig saaledes, at han paa rationel, videnskabelig Maade
søger de Love, der ligge til Grund. Dette Standpunkt havde
% footnote 1) of p. 36 runs over onto p. 37 (from „Götter als Menschen erscheinen“).
Hegel billiget overfor Mythologien\footnote{„Die Hauptsache der Mythologie ist das Werk der phantasirenden
Vernunft, die sich das Wesen zum Gegenstande macht, aber noch
kein andres Organ hat als die sinnliche Vorstellungsweise; daher die
Götter als Menschen erscheinen. Die Mythologie
% PRINTED AS IS: „kan“ for German „kann“.
kan nun studirt
werden für die Kunst u. s. w.; der denkende Geist muß aber den substantiellen
Inhalt, den Gedanken, das Philosophem, das implicite
% PRINTED AS IS: „inthalten“ for German „enthalten“.
darin inthalten ist, aufsuchen, \emph{wie man in der Natur Vernunft
sucht}“. Geschichte der Philosophie. I, p. 98.}; det var Feuerbach, der
% --- p. 37 ---
\opage{37}\setcounter{footnote}{0}%
med stor Skarpsindighed udførte dette videre og søgte at finde
de psychologiske Love, der havde ledet Religionens og navnlig
Christendommens Opstaaen og Udvikling\footnote{„Das Wesen des Christenthums“. Von Ludwig Feuerbach. Leipzig.
1841.}. Saa kom han da
til den Sætning: „Religion ist das bewußtlose Selbstbewußtsein
des Menschen“ (p. 37). Førend den menneskelige Selvbevidsthed
er ret udviklet, kan den ikke anskue sit eget almene
Væsen, Ideen som Idee, men indklæder den i en Phantasiform,
personificerer den og udstyrer den saaledes opstaaede
personlige Gud med menneskelige Attributer. Religionen er
altsaa opstaaet ved, at man har gjort til Subject, hvad der
skulde være Prædikat, og omvendt; herved er Illusionen om
en Modsætning mellem det Guddommelige og det Menneskelige
opstaaet. „Gott ist der Begriff der Gattung als Individuum“
(p. 202). „Das Bewußtsein des Menschen in seiner
Totalität ist das Bewußtsein der Trinität“ (p. 72). Altsaa
er Theologien Anthropologie. Hvad der da nu maa gjøres,
nu, da Fornuften og Selvbevidstheden er bleven myndig, er:
at gjøre det, som for Religionen er Hovedsagen, Subjectet,
til Bisagen, Prædikatet. „Das Richtige, Wahre und Gute
hat überall seinen Heiligungsgrund in sich selbst, in seiner
Qualität. Wo es Ernst mit der Ethik ist, da gilt sie eben an
und für sich selbst für eine göttliche Macht“ (p. 375 f.).

Strauß og Feuerbach sige, at de tale i Videnskabens
Navn. Men kunne de nu ogsaa virkelig gjælde for Videnskabens
rette Repræsentanter overfor Religionen? eller maa det
ikke siges at være en eensidig Opfattelse af den videnskabelige
Reflexion, at den kun skulde vidne \textit{contra}, ligesaavel som det
% --- p. 38 ---
\opage{38}\setcounter{footnote}{0}%
er eensidigt, naar Theologien kun kan høre Vidnesbyrdet \textit{pro}?
Og vilde det endelig ikke være et betænkeligt Forhold, hvori
Videnskaben kom til at staae til Humaniteten og det aandelige
Liv, naar den overfor Religionen absolut maatte være destructiv?

Hvorledes Sagen stiller sig fra Troens Side: at Troen
har sin egen Sphære, hvor den ikke erkjender Videnskabens
Suprematie (og indlod den sig først med den objektive Reflexion,
vilde den være ilde stedt: thi her er jo baade \textit{pro} og
\textit{contra}, tvende Veie!), --- det have vi allerede seet i det
Foregaaende. Det er S. Kierkegaard, der har stridt for denne,
den religieuse Side af Sagen.

Men ogsaa Videnskaben, Philosophien maa jo klare for
% „rette“: the first t's crossbar barely printed, but its stem is t-height (same as the second t, well short of an l); not a printer's error.
sig selv, hvilket dens rette Forhold til Religionen er. Fra
denne Side har Prof. R. Nielsen bearbeidet de herhen hørende
Problemer. Og det er klart, at Afgjørelsen af dette Punkt
maa hidføres ad videnskabelig Vei.

Der spørges da først: hvorfor mon Videnskabens Resultater
for Strauß og Feuerbach altid lyde \textit{contra} Aabenbaringen?
Vel af samme Grund, ifølge hvilken Theologien altid hører
et \textit{pro} af Videnskabens Mund: Fritænkerens Videnskab er en
vantro Videnskab, Theologien en troende Videnskab. Men Tro
og Vantro ere ikke videnskabelige Principer: de høre Livet,
Existentsen til.

Videnskaben spørger kun om de bestemte, objective Grunde.
„Fra Livets Standpunkt hersker derimod Subjectiviteten“; her
afgjøres da Alt ved subjective Grunde. Naar nu altsaa en
Afgjørelse af Spørgsmaal, som høre til Existentsens, det personlige
og religieuse Livs Sphære, giver sig ud for at være
videnskabelig, da er det klart, at de to Omraader forvexles,
--- hvad enten saa Afgjørelsen falder ud i den ene eller den anden
Retning. „Den straußiske Bestræbelse, at hæve Videnskaben til
Religion, kan umulig være ueensartet med den theologiske
Bestræbelse at hæve Religionen til Videnskab; thi begge Be%
% p. 38 ends mid-word at „Be-“ (hyphenated); the word continues on p. 39.
% ==== Batch 5: printed pp. 39--46 (PDF 43--50), end of the book ====
% Transcribed from the page images (300/600/900 dpi). Page 39 begins
% mid-word: p. 38 ends "Be-", so the word is Be|str\ae belser.
% --- p. 39 ---
\opage{39}\setcounter{footnote}{0}%
stræbelser gaae jo ud fra den Forudsætning, at Religion og
Videnskab ere eensartede“ (Evangelietroen og Theologien p. 82).

Der er da intet Andet for Videnskaben at gjøre end at
udtale sit \textit{non liqvet} angaaende disse Spørgsmaal og overlade
deres Afgjørelse til den Sphære, hvor de høre hjemme. Den
bliver kun komisk, naar den her kommer med sine objektive
Grunde og fylder Vægtskaalene, den ene med \textit{pro'er}, den
anden med \textit{contra'er}, --- og se! saa beroer det dog aldeles
ikke paa disse Grunde, til hvilken Side Skaalen synker.

„Men“, vil Strauß sige, „Alt i Verden hænger jo sammen
ved en Kjede af Aarsager og Virkninger; Miraklet maa da
forkastes af Videnskaben som en reen Umulighed.“ Hertil
svares fra Philosophiens Side, at saaledes maa det ganske
vist forekomme Physikeren; thi „Physikens Aarsager ere Tingsaarsager,
dens Love Nødvendighedslove; men Philosophiens
høieste Aarsag er en Idee-Aarsag, og Ideeaarsagen en teleologisk
Selvaarsag. Idet den absolute Selvaarsag lægger Nødvendighedslovene
til Grund for Frihedens Lov, sees fra Idee-Erkjendelsens
Høide et Dyb af Muligheder, som ingen menneskelig
Viden er istand til at udmaale. Istedetfor at phantasere
over det Ubestemmelige og gruble over det Uudgrundelige,
holder Philosophien her fast ved den væsentlige Opgave:
at sikkre Videnskabens Ret ved at angive Videns Grændse“%
% fn 1) p. 39
\footnote{Grundideernes Logik. I, p. XXIV.}.
Philosophien kunde altsaa umuligt hævde Aandens og det
Ideelles Gyldighed i Tilværelsen, hvis den vilde \emph{bryde} med
Religionen, hvor Aandens absolute Almagt hævdes. Gjælder
dette paa det intellectuelle Omraade, da gjælder det end mere
paa det ethiske. Thi det forholder sig ikke, som Feuerbach
mener, at „das Richtige, Wahre und Gute“ vilde kunne holde
sig i det virkelige Liv, naar det ikke støttedes af Religionen.
Det Godes Idee vilde snart blive bortelimineret af den menneske\-%
% --- p. 40 ---
\opage{40}\setcounter{footnote}{0}%
lige Reflexion%
% fn 1) p. 40
\footnote{Dette udvikles udførligt i „Forelæsninger over philos. Propædeutik
1860--61“. p. 416--423.}. Saa vist derfor, som Philosophien skal
have en Ethik, hvor „Livstankerne mødes med de videnskabelige
Grundtanker“, saa vist er det, at den ikke tør erklære Evangeliet
for en Mythe (Grundideernes Logik XXVII).

Men begribe og begrunde Evangeliet kan Philosophien
ligesaa lidt. „Idet Videnskaben erklærer, at Objektiviteten er
Sandheden, fraskriver den sig med det Samme al Indflydelse
paa Troen“.

Derved er imidlertid ikke udelukket, at Religionen og
overhovedet det personlige Liv kan blive Gjenstand for Videnskab.
Men dette kan da kun skee derved, at det Subjective
„objectiveres i Reflexionen“, og den religieuse Anskuelse saaledes
haves „paa anden Haand“, „ikke paa religieus Maade“, men
„gjennem Tankens og Phantasiens Medium“ (cfr. Martensens
Dogmatik § 8 Anm.). Her maa man da holde fast ved den
kategoriske „Forskjel mellem det i videre og det i snevrere Forstand
Videnskabelige“ (R. Nielsens philos. Propædeutik p. 78).
Kun i videre Forstand videnskabelig er enhver Behandling af
en Gjenstand, som ligger udenfor den Sphære, hvor Videnskaben
hersker med absolut \textit{dominium}; altsaa enhver Fremstilling
af det subjektive, personlige Liv. Det Videnskabelige
ligger i den objektive Behandling. En saadan „Subjektivitetsvidenskab“
kan da have de religieuse Sandheder til Gjenstand,
ikke for at begrunde dem, men for at „belyse dem saa alsidigt
som muligt“, idet den „opfatter Troesidealerne som Livsidealer
og i Formaalet finder Maalestokken for en Bedømmelse af
Midlerne“ (Grundideernes Logik XXVI).

Forholder det sig nu saaledes, at der er en qvalitativ
Forskjel mellem Tro og Viden, at de høre til ganske forskjellige
Sphærer, saa er herved Muligheden givet til deres
Forsoning. Hvor, paa hvilket Omraade, i hvilken Sphære
% --- p. 41 ---
\opage{41}\setcounter{footnote}{0}%
skal denne Forsoning gaae for sig? I Videns Sphære? Men
saa er Viden selv baade Part og Dommer; Mediationen vil
derfor blive partisk, Eenheden eensidig (jvnfr. Uvidensk. Efterskr.
p. 286 f.). I Troens Sphære? Her kan der svares Ja,
naar man ved „Troens Sphære“ forstaaer den personlige
Existerens Sphære i det Hele. Det er den existerende Subjectivitet,
som har et Forhold baade til Viden og til Troen, i
hvis Bevidsthed Striden føres. „Striden er psychologisk; det
er den menneskelige Skrøbelighed, der hverken har Resignation
nok til at forsone sig med Videnskabens Objective, eller Inderlighed
nok til at fastholde Troen“ (Philos. Propæd. p. 77).

Hvis Ueensartetheden mellem Tro og Viden ikke var
absolut, men kun relativ, saa at begge hørte til samme Sphære,
hvorledes vilde de da kunne forenes i een Bevidsthed uden
Modsigelse? Naar Troen siger Et, Videnskaben et ganske
Andet og Modsat, hvorledes vilde man da kunne forsone disse
stridende Paastande? Erkjendes det derimod, at Et er existentiel
og religieus, et Andet videnskabelig Sandhed, da kan det
% PRINTED AS IS: antiqua guillemets, both printed as » (opening and closing)
gjælde: »\textit{Suum cuiqve!}» Da er det klart, at den psychologiske
Forsoning \emph{kan} (thi om den \emph{virkelig} kommer istand
beroer jo paa den enkelte existerende Subjectivitet) finde Sted,
saaledes nemlig, at Troen sættes øverst, medens Videnskaben
dog beholdes som en relativ Opgave; thi det Omvendte er
ikke muligt%
% fn 1) p. 41
\footnote{Altsaa netop \emph{fordi} Tro og Viden ere absolut ueensartede, er deres
psychologiske Forsoning mulig. Det er en stor Misforstaaelse istedetfor
dette „fordi“ at sætte et „\emph{skjøndt}“ (se Hr. \textit{cand. mag.} Brandes's
Artikel i „Fædrelandet“ den 13de Januar). Saa kan man sagtens
finde „Dunkeltheder“!}.

Forsoningen mellem Tro og Viden slaaes altsaa i den
% stray mark (spacing material?) between "Philosophie" and "ikke" -- not a hyphen
% "tvert-": the r is overlaid by an ink blot, letter form r; reading tvertimod
Nielsenske Philosophie ikke dogmatisk fast. Enhver faaer tvertimod
en Anviisning og Opfordring til praktisk og ethisk at
fuldbyrde Forsoningen under sin egen personlige Livsudvikling.
Hvad Philosophien kan lære os, er at holde Principerne fast,
% --- p. 42 ---
\opage{42}\setcounter{footnote}{0}%
ikke at forvanske og sammenblande dem. Synes det saa, som
om Striden ikke kunde bilægges, som fik Kampen aldrig Ende, ---
da gjælder Ørsteds Ord: „Lader os give Sandheden Æren!
med den er det Gode uopløseligt forbundet. Den fulde Sandhed
bringer selv sin Trøst med sig.“ Og Trøsten er den, at man
ikke har fornegtet Idealet.

Det maa her tillige erindres, at Videnskaben kun er en
enkelt, om end en saare betydningsfuld, Side af det menneskelige
Liv. Er Talen derfor om Forholdet mellem Christendom
og Videnskab, da er dette kun en særegen Form af
Spørgsmaalet om Forholdet mellem Christendom og Humanitet
i det Hele. Videnskaben staaer ved Siden af Kunsten, Familielivet
og det borgerlige Liv, kort: Alt, hvad der hører til det
reent Menneskelige, det Relative overfor det, der i Christendommen
stilles som det høieste Formaal: „Guds Rige“.

Ogsaa i disse andre Retninger er en Sammenblanding
af det Christelige og det Humane mulig, ja den har fundet
Sted, og den finder sikkert Sted endnu. Ogsaa her gjælder
det ikke at blive for tidligt færdig med „Forsoningen“, men
at fatte Modsætningerne saa skarpt som muligt i Øie --- og
saa at tage sit Parti.

I den hellige Skrift finder man ingen ligefremme Oplysninger
om Forholdet mellem Tro og Viden, mellem Christendom
og Humanitet. Dette har sin Grund deri, at Christus
og Apostlene stode paa en saadan Existentsens Høide, saa
absolut paa det religieuse Standpunkt, at det --- i Forhold
hertil relative --- humane Ideerige blev noget aldeles Forsvindende.
Men dette gjælder kun for den, der staaer paa
Inspirationens Standpunkt, hvor den hellige Aand virker ubetinget
og umiddelbart, uden at gaae ind paa og træde i Vexelvirkning
med de givne Betingelser (cfr. Apostlenes Gjerninger
2, 3--4).

I den følgende Tid udviklede sig den menneskelige Cultur
i Vexelvirkning med, ja en lang Tid under Formynderskab af
% --- p. 43 ---
\opage{43}\setcounter{footnote}{0}%
Religionen. Men Bevidstheden om en Modsætning mellem
begge kommer bestandigt igjen. Bliver Modsætningen nu en
Dualisme, saa er der kun tvende Veie mulige: Asketens eller
Fritænkerens. Den ene bryder med Humaniteten, den anden
med Christendommen.

Men lad os see nærmere til: er Asketens Standpunkt
Evangelietroens, den oprindelige Christendoms? Man behøver
ikke at have læst meget i det nye Testamente for at vide, at
et voldsomt og smerteligt Brud mellem Christendom og Humanitet,
som det, vi finde i Asketens Liv, findes ikke hos
Apostelen. Denne forener i Aandens Kraft og Høihed de afgjørende
Modsætninger, uden at man mærker nogen Smerte,
nogen urolig Uklarhed hos ham; hvorimod Asketen vaander sig
som under en tung Byrde.

Altsaa: „Respect for Middelalderens Klosterbevægelse!“
(Joh. Climacus) Respect for den, der kan bære det Ubetingede
ubetinget i denne Betingelsernes Verden! Men det asketiske
Forbillede tør ikke fordunkle det evangeliske. Evangeliet har
hverken lyst Interdict over Videnskaben eller Staten eller noget
almeen-menneskeligt Vilkaar.

Heraf sluttes, at det maa være muligt at forene det
Christelige og det Humane, at en Dualisme umuligt kan være
det rette Forhold.

Er Christendommen virkelig til endnu, er den ikke blevet
overvundet af den fremskridende Cultur, saa maa det være muligt
at fastholde det evangeliske Ideal midt i den humane Verden.
% stray mark between "lever" and "paa" -- not punctuation
Men ligesaa vist er det, at saa maa man erkjende, at man
lever paa andre Aandsvilkaar end de apostoliske og nytestamentlige.

Thi fremfor Alt: ingen Sammenblanding. Og en
Sammenblanding finder Sted, hvor man ikke agter den
væsentlige og afgjørende Modsætning. Saa faaer man hverken
sand Christendom eller sand Humanitet. ---

Naar Forsoningen af Tro og Viden skal være psychologisk,
existentiel, saa sees, hvor lidet træffende den Indvending er, at
% --- p. 44 ---
\opage{44}\setcounter{footnote}{0}%
man ved at henvise Tro og Viden til qvalitativt forskjellige
Sphærer sikkrer sig sit Troesstandpunkt mod alle Fritænkeres
Indsigelser: „hvad Risico er der i at troe „mod Begrebet“,
naar vi iforveien have sikkret os den Indsigt, at Begrebet aldeles
ingen Gyldighed har i Troens Sphære?“ („Dogmatiske
Oplysninger“ p. 74). Thi fordi det er erkjendt \emph{videnskabeligt},
at Viden ingen Myndighed har paa Troens Omraade, mon
det derfor bliver lettere at \emph{troe}? Tvertimod, den Indsigt, som
her er vunden i \emph{Viden}, maa helt forfra og paa andre Vilkaar
vindes igjen, naar det gjælder at \emph{troe}. Thi „ikke kan det
opstilles som et nyt Dogma om Dogmet, at Dogmet er et
Tankekors. Kun ved Forsøget paa at løfte Dogmet gjør
Reflexionskraften denne Opdagelse“%
% fn 1) p. 44
% "og --- Viden": the dash stands in the print
\footnote{Stilling: „Om den indbildte Forsoning af Tro og --- Viden“.
p. 55 Anm.}. Og hvad det angaaer,
der tales om „de eenfoldige Troende, der uden en saadan
Forberedelse sætte sig paa Troens Standpunkt“, da maa bemærkes,
først: at det, som sagt, er en Misforstaaelse, naar
der tales om „en saadan Forberedelse“, og dernæst: at her
gjælder den samme „eenfoldige“ Forskjel mellem „den Vise“
og „den Eenfoldige“ som ved al ethisk Erkjendelse: „den Eenfoldige
veed det; den Vise veed af, at han veed det“ (Joh.
Climacus). Iøvrigt er det vel neppe fra den spekulative
Theologies Side, at der skulde tales for Meget om „de eenfoldigt
Troende“.

Medens man her har søgt at fremstille Anskuelsen om
en Sphæreforskjel mellem Tro og Viden som en ny „Scholasticisme“,
saa er der Andre, der beskylde den for videnskabelig
„Skepticisme“. Det „frapperer“ Hr. Larsen (i „Dagbladet“
af 19de Januar)%
% fn 2) p. 44
\footnote{Af denne Artikel sees iøvrigt, at har Hr. Larsen ikke lært Andet af
den tidligere Strid om denne Sag, saa har han idetmindste lært at
lyse i Ban.}, han „kan paa ingen Maade faae det i
sit Hoved“, hvorledes man kan faae en Videnskab, der erkjender
% --- p. 45 ---
\opage{45}\setcounter{footnote}{0}%
„evige Naturlove“, uden at man ligefrem erklærer Underet for
umuligt. Han angiver da „tre mulige Standpunkter“ overfor
Underet, men til liden Nytte for Læseren; thi man faaer ikke
at vide, om alle disse tre Standpunkter ere videnskabelige, eller
om et eller to af dem er det, eller om maaskee slet intet af
dem er det. Og herpaa var det dog, det Hele kom an. Herom
skulde han have raadspurgt sin „kyndige Bygningskommission“.
Sammesteds kunde han maaskee ogsaa skaffe mig at vide, hvad
der „frapperer“ mig, hvad jeg ikke „kan faae i mit Hoved“,
nemlig: hvorledes bærer man sig ad med at drage Underet
ind i Videnskaben og ved Siden heraf at holde paa Naturvidenskaben
med de „evige Naturlove“?

Det samme Spørgsmaal, som her behandles saa ubehændigt,
kommer paa en langt klarere Maade frem i Hr. G.
Brandes's Artikler i „Fædrelandet“ af 13de og 20de Januar.
Jeg forstaaer Spørgsmaalet saaledes: hvorledes kan Philosophien
paa sit eget Omraade videnskabeligt retfærdiggjøre det
Hensyn, den tager til en den fremmed Sphære? --- Hermed
er da Diskussionen ført bort fra Grændsestriden mellem Philosophie
og Theologie til det reent philosophiske Omraade%
% fn 1) p. 45
\footnote{Hr. Brandes staaer i et særeget Forhold til Hr. Larsen. Ved første
Øiekast skulde man troe, at han forsvarede Hr. Larsen; men saa erklærer
han dennes Forsvar for Theologien for betydningsløst og
tilskriver ham kun Fortjeneste som Angriber; og se, saa er det Punkt,
hvor Hr. Brandes mener, at Angrebet maa skee, et ganske andet
end det, som Hr. L. har angrebet. Følgelig beholder Hr. L. kun
den Fortjeneste tilbage, at han overhovedet har angrebet Prof. Nielsen.}.
Forsaavidt vedkommer den egentlig ikke nærværende Afhandling.
Imidlertid vil jeg bemærke Følgende.

Hvad der forlanges, kan kun være den explicite Udvikling
af, hvad det vil sige, at „Philosophiens høieste Aarsag“ er
„den absolute Selvaarsag“ (Grundideernes Logik XXIV).
Denne Udvikling kan, efter Hr. Brandes's Mening, først skee
ved Overgangen fra Videns Kreds til Magtens og, i endnu
% --- p. 46 ---
\opage{46}\setcounter{footnote}{0}%
concretere Form, ved Overgangen fra Logiken til Naturphilosophien.
Men hvis Hr. Brandes virkelig slet ikke har fundet
Momenter til sit Problems Løsning i det hidtil Udkomne af
„Grundideernes Logik“, saa skulde han ikke have indskrænket
sig til at ønske Oplysning og Fuldstændiggjørelse, men saa
maatte han have fundet Feil allerede i Udviklingen af „Subjectivitetens
Logik“ og „Functionernes Logik“. Thi hiint Savn
forudsætter nødvendigt denne Feil, hvis ellers det Udkomne af
„Grundideernes Logik“ virkelig er et Hele for sig og ikke et
Brudstykke.

Derhos undrer det mig, at Hr. Brandes ikke har taget
Hensyn til den Betydning, der i Indledningen til „Grundideernes
Logik“ (p. XXVII) tillægges Ethiken. Thi skal der
gives en philosophisk Ethik, da kan denne ikke staae i Modsigelse
med Logiken og Naturphilosophien. Og idet Ethiken
danner en Overgang mellem Philosophie og Religion, er herved
Veien anviist, ad hvilken man, uden at „gjøre Videnskaben
religieus“ eller „skrive en ny Theologie“, kan respectere det
Religieuse ved at erkjende Videnskabens Grændse. ---

Her kan være Meget at gjøre endnu. Det er heller ikke
forjættet, at alle Spørgsmaal skulle kunne finde deres absolute
Løsning. „Det forholder sig ikke med vor Videns Fremskridt
saaledes, at de resterende Problemer skulde blive færre, lettere
og ubetydeligere, jo mere vi nærme os Idealet; tvertimod! jo
klarere Idealet sees, desto høiere blive Problemerne og desto
dybere gaaer Undersøgelsernes Strøm“%
% fn 1) p. 46
\footnote{Grundideernes Logik p. 253.}. --- Men hvad der
kan opnaaes er, at der gjøres væsentlige Fremskridt, at man
bliver rigtigere stillet. Og heri ligger efter min Overbeviisning
den store Betydning af den her fremstillede religionsphilosophiske
Anskuelse, hvis historiske Forudsætninger, Udvikling og
væsentlige Indhold jeg i denne lille Afhandling har stræbt at
tydeliggjøre.

% Closing rule, centred, below the footnote on p. 46 (probably a double
% rule; the scan cannot resolve it). No date, signature or printer's line.
% PDF 51 = blank end-paper (library pencil note only), PDF 52 = back board.
\begin{center}\rule{0.2\textwidth}{0.8pt}\end{center}

\end{document}
